% =====================================================================
% Chapter: Mass as a Momentum in a Hidden Dimension
%          Galilean Spacetime as a Null Reduction
%
% Requires: amsmath, amssymb, amsthm, booktabs
% All theorem-like environments and labels carry the prefix "nr"
% (null reduction) so they do not collide with the rest of the book.
% =====================================================================

\theoremstyle{plain}
\newtheorem{nrprop}{Proposition}[chapter]
\theoremstyle{definition}
\newtheorem{nrdef}[nrprop]{Definition}
\theoremstyle{remark}
\newtheorem{nrrem}[nrprop]{Remark}
\theoremstyle{plain}

\chapter[Mass as a Momentum in a Hidden Dimension]{Mass as a Momentum in a Hidden Dimension:\\ Galilean Spacetime as a Null Reduction}
\label{ch:nr}

%---------------------------------------------------------------------
\section{Introduction}
%---------------------------------------------------------------------

Mass is the most familiar and, in a sense, the least explained quantity of
mechanics. In Newtonian practice it is a positive number attached to a body.
Once one asks what role it plays in the \emph{symmetry} structure of the
theory, two very different answers appear, depending on which kinematics one
uses.

In Galilean mechanics, mass is a \emph{central charge}: it appears on the
right-hand side of the Poisson bracket of the boost charge with the momentum,
$\{K_i,P_j\}=m\,\delta_{ij}$, and it cannot be removed by any redefinition of
the generators. In relativistic mechanics there is no such term. The
Poincar\'e algebra has no non-trivial central extension, and mass is instead
the value of a \emph{Casimir invariant}, the squared length of the
four-momentum. The two statements look unrelated, and the textbook way of
connecting them (the limit $c\to\infty$) is a limit of algebras, which hides
what is going on geometrically.

This chapter shows that both statements are two faces of one construction.
Galilean spacetime is the space of orbits of a \emph{null} translation in a
Minkowski space with one more spatial dimension, and mass is the momentum
conjugate to the null coordinate. Nothing is approximated, no parameter is sent
to infinity, and the central extension of the Galilei algebra becomes an
ordinary translation generator of a larger, perfectly honest, geometry.

\paragraph{What will be established.}
\begin{enumerate}
\item In Galilean mechanics the mass enters the Poisson algebra of conserved
      charges as a non-trivial central term. It cannot be removed, and it is
      the algebraic trace of the fact that the Lagrangian is only
      \emph{quasi}-invariant under boosts (Section~\ref{sec:nr-galilei}).
\item Galilean spacetime is the quotient of $\mathbb{R}^{d+1,1}$ by the flow of
      a null parallel vector field $\xi$; the quotient inherits exactly the
      Galilean structure, absolute time plus Euclidean slices, and nothing more
      (Section~\ref{sec:nr-geometry}).
\item Free Newtonian motion is the projection of null geodesics. Mass is the
      conserved momentum $p_s$ along $\xi$, and the Newtonian action equals
      $-m$ times the displacement of the null coordinate
      (Section~\ref{sec:nr-dynamics}).
\item The Bargmann group, the central extension of the Galilei group, is the
      stabiliser of $\xi$ in the Poincar\'e group of the ambient space. The
      cocycle is realised as a shift of the extra coordinate
      (Section~\ref{sec:nr-symmetry}).
\item Newtonian gravity, and the universality of free fall, follow from adding
      one term to the ambient metric (Section~\ref{sec:nr-gravity}).
\item The same mathematics is already present in relativistic light-front
      dynamics, where the Casimir of the Bargmann group is the Poincar\'e mass
      squared (Section~\ref{sec:nr-lightfront}), and it differs in kind from
      the $c\to\infty$ contraction (Section~\ref{sec:nr-limit}).
\item It is used in practice to couple Galilean systems to background fields,
      in the quantum Hall effect, and in non-relativistic holography, where
      the Schr\"odinger spacetime is a null-reduction geometry
      (Section~\ref{sec:nr-applications}).
\end{enumerate}

\paragraph{Conventions.}
Latin indices $i,j,\dots=1,\dots,d$ label space ($d=3$ in the physical case;
we keep $d$ general because the counting is instructive). Greek indices
label the $d+2$ coordinates $(x^i,t,s)$ of the ambient space. Poisson brackets
are $\{A,B\}=\partial_{q}A\,\partial_{p}B-\partial_{p}A\,\partial_{q}B$.
The reduced Planck constant appears only in Section~\ref{sec:nr-quantum}, where
$\hbar=1$. The ambient geometry contains no speed of light; the only place
where $c$ appears is Section~\ref{sec:nr-limit}.

%---------------------------------------------------------------------
\section{Mass as a central charge in Galilean mechanics}
\label{sec:nr-galilei}
%---------------------------------------------------------------------

\subsection{Boosts and the quasi-invariance of the Lagrangian}

Take a free particle, $L=\tfrac12 m\dot{\mathbf{x}}^2$. Under a Galilean boost
with constant velocity $\mathbf{u}$,
$\mathbf{x}\mapsto\mathbf{x}'=\mathbf{x}+\mathbf{u}t$, $t'=t$, the velocity
shifts by $\mathbf{u}$ and
\begin{equation}
  L(\dot{\mathbf{x}}')
  = \tfrac12 m(\dot{\mathbf{x}}+\mathbf{u})^2
  = L + m\,\mathbf{u}\cdot\dot{\mathbf{x}} + \tfrac12 m u^2
  = L + \frac{\mathrm{d}}{\mathrm{d}t}\Bigl(m\,\mathbf{u}\cdot\mathbf{x}
      + \tfrac12 m u^2 t\Bigr).
  \label{nr:quasi}
\end{equation}
The Lagrangian is not invariant, but it changes by a total time derivative
$\mathrm{d}F_{\mathbf u}/\mathrm{d}t$, so the action changes only by boundary
terms and the equations of motion are invariant. This is called
\emph{quasi-invariance}. The function
$F_{\mathbf u}=m\,\mathbf{u}\cdot\mathbf{x}+\tfrac12 m u^2 t$
is proportional to $m$, and we will see that it carries the whole of the
central extension.

Noether's theorem extends to quasi-symmetries. If under
$\delta\mathbf{x}=\boldsymbol{\varepsilon}\,t$ one has
$\delta L=\mathrm{d}F/\mathrm{d}t$ with $F=m\,\boldsymbol{\varepsilon}\cdot\mathbf{x}$
(the part of \eqref{nr:quasi} linear in $\mathbf u$), the conserved quantity is
$Q=\mathbf{p}\cdot\delta\mathbf{x}-F=\boldsymbol{\varepsilon}\cdot(t\,\mathbf{p}-m\mathbf{x})$.
Up to sign, the conserved boost charge is
\begin{equation}
  \mathbf{K}=m\,\mathbf{x}-t\,\mathbf{p},\qquad
  \frac{\mathrm{d}\mathbf{K}}{\mathrm{d}t}
  =\frac{\partial\mathbf{K}}{\partial t}+\{\mathbf{K},H\}
  =-\mathbf{p}+\mathbf{p}=0 .
  \label{nr:Kcons}
\end{equation}
Its conservation is the statement $\mathbf{x}(t)=\mathbf{x}_0+(\mathbf{p}/m)t$:
the free motion with $\mathbf{p}=m\mathbf{v}$. Note where the mass came from.
The term $-t\mathbf{p}$ comes from the transformation of the coordinates, but
the term $m\mathbf{x}$ comes \emph{entirely} from the total derivative $F$. The
mass sits in the boost charge only through the amount by which the Lagrangian
fails to be invariant.

\subsection{The Poisson algebra of conserved charges}

The ten Noether charges of a free particle in $d=3$ (in general
$\tfrac12 d(d-1)+2d+1$ of them) are
\begin{equation}
  P_i=p_i,\quad H=\frac{\mathbf{p}^2}{2m},\quad
  J_{ij}=x_ip_j-x_jp_i,\quad K_i=m\,x_i-t\,p_i ,
  \label{nr:charges}
\end{equation}
and, trivially, the mass $M=m$ as a number that commutes with everything. With
$\{x_i,p_j\}=\delta_{ij}$ one finds directly
\begin{align}
  \{K_i,P_j\} &= \{m x_i,p_j\} = m\,\delta_{ij}, \label{nr:KP}\\
  \{K_i,H\}   &= \Bigl\{m x_i,\frac{\mathbf{p}^2}{2m}\Bigr\} = p_i = P_i, \label{nr:KH}\\
  \{K_i,K_j\} &= 0, \qquad \{P_i,H\}=0,\qquad \{M,\cdot\,\}=0, \label{nr:KK}
\end{align}
while the rotations $J_{ij}$ act on the indices of $P_i$ and $K_i$ as on
vectors. Equations \eqref{nr:KP}--\eqref{nr:KK} are the \emph{Bargmann}
algebra. If one removes $M$ and sets $\{K_i,P_j\}=0$ one obtains the Galilei
Lie algebra, whose abstract bracket is
\begin{equation}
  [K_i,P_j]=0,\qquad [K_i,H]=P_i ,
  \label{nr:galilei}
\end{equation}
with the same rotation relations. The map $X\mapsto Q_X$ from the Galilei
algebra to the charges \eqref{nr:charges} is therefore \emph{not} a Lie algebra
homomorphism: $Q_{[K_i,P_j]}=Q_0=0$, but $\{Q_{K_i},Q_{P_j}\}=m\delta_{ij}\neq0$.

\subsection{The term cannot be removed}

The defect of the homomorphism property has a precise name. For any
action of a Lie algebra $\mathfrak{g}$ on a phase space with charges $Q_X$,
\begin{equation}
  \Theta(X,Y)\;:=\;\{Q_X,Q_Y\}-Q_{[X,Y]}
  \label{nr:theta}
\end{equation}
is a constant (it Poisson-commutes with everything) whenever the
$\mathfrak{g}$-action is Hamiltonian, and it is a \emph{2-cocycle}: it is
antisymmetric and satisfies the Jacobi-compatibility condition
$\Theta([X,Y],Z)+\text{cyclic}=0$. This is Souriau's formulation: the moment
map of the action is not equivariant, and the failure of equivariance is the
cocycle $\Theta$. For the free particle the only non-zero component is
$\Theta(K_i,P_j)=m\,\delta_{ij}$.

One may try to absorb a cocycle by adding constants to the charges,
$Q_X\mapsto Q_X+\lambda(X)$ with $\lambda$ a linear functional. This changes
$\Theta$ by $\lambda([X,Y])$, a \emph{coboundary}.

\begin{nrprop}\label{nr:prop-nontrivial}
The cocycle $\Theta(K_i,P_j)=m\,\delta_{ij}$ of the Galilei algebra is not a
coboundary for $m\neq0$.
\end{nrprop}
\begin{proof}
A coboundary has the form $\Theta(X,Y)=\lambda([X,Y])$. By \eqref{nr:galilei},
$[K_i,P_j]=0$, so every coboundary vanishes on $(K_i,P_j)$, whereas
$\Theta(K_i,P_j)=m\delta_{ij}\neq0$.
\end{proof}

So no shift of the charges can remove the mass from the algebra. In the
language of cohomology, the second cohomology group of the Galilei algebra for
$d=3$ is one-dimensional, and mass is its generator (Bargmann 1954,
L\'evy-Leblond 1963).

\begin{nrrem}[Why the relativistic case is different]
In the Poincar\'e algebra the corresponding bracket is not zero:
$[K_i,P_j]=\delta_{ij}H/c^2$ (up to a convention-dependent sign). A candidate
central term $\kappa\,\delta_{ij}$ in this bracket is then a coboundary: it is
absorbed by the constant shift $H\mapsto H+c^2\kappa$. In fact all
second cohomology of the $3{+}1$ Poincar\'e algebra vanishes, so there is no
central extension at all. The mass cannot appear as a central charge in
relativity; it appears as the value of a Casimir function (Section~\ref{sec:nr-lightfront}).
This is a precise form of the fact that in relativity there is no separately
conserved, additive ``mass'' to be split off from the energy: the generator $H$
of time translations already contains the rest energy, and the algebra leaves no
room for a second, independent central generator.
\end{nrrem}

\begin{nrrem}[A second extension in $d=2$]
For $d=2$ spatial dimensions the Galilei algebra admits a second, ``exotic''
central extension, in which the boosts do not commute,
$\{K_1,K_2\}=\theta$ (L\'evy-Leblond; see also Duval and Horv\'athy). It plays no
role below.
\end{nrrem}

\subsection{Composite systems: additivity and internal energy}

For $N$ particles with masses $m_a$ the same computation gives
\begin{equation}
  \mathbf{K}=\sum_a m_a\mathbf{x}_a-t\,\mathbf{P},\qquad
  \{K_i,P_j\}=M\,\delta_{ij},\quad M=\sum_a m_a .
  \label{nr:Mtotal}
\end{equation}
The central charge of a composite system is the \emph{sum} of the central
charges of its constituents, and it does not depend on how they interact: for
any interaction potential depending only on relative positions, $M$ remains the
sum of the bare masses. Writing the centre of mass as
$\mathbf{X}=\sum_a m_a\mathbf{x}_a/M$, one has $\mathbf{K}=M(\mathbf{X}-t\mathbf{V})$
with $\mathbf{V}=\mathbf{P}/M$, and the energy splits as
\begin{equation}
  H=\frac{\mathbf{P}^2}{2M}+H_{\rm int},
  \label{nr:Hint}
\end{equation}
where $H_{\rm int}$ depends only on relative variables. One checks that
$\{X_i,H_{\rm int}\}=0$ and $\{P_i,H_{\rm int}\}=0$, hence $H_{\rm int}$
commutes with all of $\mathbf{K}$ and $\mathbf{P}$.

The Bargmann algebra has the following independent Casimir functions for $d=3$:
\begin{equation}
  M,\qquad
  C=2MH-\mathbf{P}^2=2M\,H_{\rm int},\qquad
  \mathbf{S}^2,\quad \mathbf{S}=\mathbf{J}-\frac1M\,\mathbf{K}\times\mathbf{P},
  \label{nr:Casimirs}
\end{equation}
where $J_k=\tfrac12\epsilon_{kij}J_{ij}$. For instance, from \eqref{nr:KP} and
\eqref{nr:KH},
\[
  \{K_i,2MH-\mathbf{P}^2\}=2M P_i-2P_j\,M\delta_{ij}=0 .
\]
Similarly $\{K_i,S_j\}=\epsilon_{ijk}K_k-\epsilon_{ijk}K_k=0$ and
$\{P_i,S_j\}=0$, so $\mathbf{S}$ (the internal angular momentum) commutes with
boosts and translations and $\mathbf{S}^2$ is a Casimir. In Souriau's
classification, the elementary Galilean systems (the coadjoint orbits of the
Bargmann group) are labelled by three numbers: the mass $M$, the internal
energy $U=H_{\rm int}$, and the spin $|\mathbf{S}|$.

\subsection{The quantum shadow: Bargmann's superselection rule}

The central term has a classical meaning and a quantum consequence. Let
$T_{\mathbf a}$ be a spatial translation and $B_{\mathbf v}$ a boost. As
transformations of Galilean spacetime they commute, since
$\mathbf{x}\mapsto\mathbf{x}+\mathbf{a}+\mathbf{v}t$ is the same in either
order. In a quantum theory in which \eqref{nr:KP} is realised, the
corresponding unitary operators satisfy
$B_{\mathbf v}T_{\mathbf a}B_{\mathbf v}^{-1}T_{\mathbf a}^{-1}=e^{\pm i m\,\mathbf{v}\cdot\mathbf{a}}$
(with $\hbar=1$): a pure phase, allowed because states are rays. This phase
depends on $m$. If a state were a superposition of components with different
masses $m_1\ne m_2$, the group commutator, which acts as the identity on
spacetime, would multiply the two components by different phases and so map
the state to a different ray. Hence superpositions of states with different
masses cannot be physical states: mass is a superselection label. We will see
in Section~\ref{sec:nr-quantum} where this phase comes from geometrically.

%---------------------------------------------------------------------
\section{The geometry of the null reduction}
\label{sec:nr-geometry}
%---------------------------------------------------------------------

\subsection{The ambient space}

Let $\tilde M=\mathbb{R}^{d+2}$ carry coordinates $(x^1,\dots,x^d,t,s)$ and the
flat metric
\begin{equation}
  g=\delta_{ij}\,\mathrm{d}x^i\mathrm{d}x^j+2\,\mathrm{d}t\,\mathrm{d}s .
  \label{nr:metric}
\end{equation}
This is Minkowski space $\mathbb{R}^{d+1,1}$ (signature $d+1$ plus signs and
one minus sign; $d=3$ gives five dimensions) written in null coordinates for a
plane. Explicitly, with ordinary Minkowski coordinates $X^0$ and $X^{d+1}$,
\begin{equation}
  t=\frac{X^{d+1}+X^0}{\sqrt2},\qquad s=\frac{X^{d+1}-X^0}{\sqrt2},\qquad
  -(\mathrm{d}X^0)^2+(\mathrm{d}X^{d+1})^2=2\,\mathrm{d}t\,\mathrm{d}s .
  \label{nr:nullcoords}
\end{equation}
The inverse metric has $g^{ij}=\delta^{ij}$, $g^{ts}=g^{st}=1$ and all other
components zero. Dimensionally, $t$ is a time and $s$ has the dimensions
of $(\text{length})^2/\text{time}$, so that $\mathrm{d}t\,\mathrm{d}s$ is a
length squared. No speed of light has entered.

\subsection{A null Killing vector and the quotient}

Consider the vector field
\begin{equation}
  \xi=\frac{\partial}{\partial s}.
  \label{nr:xi}
\end{equation}
Its properties are the following.
\begin{itemize}
\item It is \emph{null}: $g(\xi,\xi)=g_{ss}=0$.
\item It is a \emph{Killing} vector: the components of $g$ do not depend on $s$,
      so $\mathcal{L}_\xi g=0$.
\item It is \emph{parallel}: all Christoffel symbols vanish in these
      coordinates, so $\nabla\xi=0$.
\item Its dual one-form is $\xi^\flat=g(\xi,\cdot)=\mathrm{d}t$.
\end{itemize}
The flow of $\xi$ moves a point along the $s$-direction at fixed $(\mathbf{x},t)$.
The orbit space is
\begin{equation}
  M=\tilde M/\mathbb{R}\;\cong\;\mathbb{R}^{d+1}\ni(\mathbf{x},t),\qquad
  \pi:\tilde M\to M,\quad \pi(\mathbf{x},t,s)=(\mathbf{x},t).
  \label{nr:quotient}
\end{equation}
We call $M$ the \emph{base}, or Galilean spacetime, and $\pi$ the null
reduction. The question is which parts of the ambient geometry survive the
quotient.

\subsection{The induced Galilean structure}

\begin{nrprop}\label{nr:prop-induced}
The following structures on $\tilde M$ descend to $M$:
\begin{enumerate}
\item the \emph{clock form} $\tau=\xi^\flat=\mathrm{d}t$;
\item a positive-definite metric $h$ on the vectors annihilated by $\tau$,
      namely $h=\delta_{ij}\mathrm{d}x^i\mathrm{d}x^j$ on each slice $t=\text{const}$.
\end{enumerate}
The spacetime metric $g$ itself does not descend.
\end{nrprop}
\begin{proof}
(1) $\tau(\xi)=g(\xi,\xi)=0$ and $\mathcal{L}_\xi\tau=0$, so $\tau$ is a basic
form, i.e.\ the pull-back of a form on $M$.
(2) Let $V,W$ be tangent vectors on $M$ with $\tau(V)=\tau(W)=0$, and let
$\tilde V,\tilde W$ be lifts to $\tilde M$. The lifts are determined only up to
adding multiples of $\xi$. Then
\[
  g(\tilde V+a\xi,\tilde W+b\xi)
  =g(\tilde V,\tilde W)+a\,\tau(W)+b\,\tau(V)+ab\,g(\xi,\xi)
  =g(\tilde V,\tilde W),
\]
so $h(V,W):=g(\tilde V,\tilde W)$ is well defined. For a general vector, say
$V=\partial_t$, the term $a\,\tau(W)$ does not vanish and $g(\tilde V,\tilde W)$
depends on the choice of lift, so $g$ does not descend.
\end{proof}

Equivalently, in terms of covectors: the inverse metric restricted to covectors
annihilating $\xi$ (those with $p_s=0$) is $g^{-1}(p,p)=\mathbf{p}^2$, a
degenerate ``cometric'' $h^{\mu\nu}=\delta^{ij}$ on $M$ of rank $d$ whose kernel
is spanned by $\tau$.

The pair $(\tau,h)$ is precisely Galilean (Newton--Cartan) structure on flat
spacetime: there is an absolute time function $t$, the notion of
simultaneity is given by $\tau=0$, and each instantaneous space carries a
Euclidean geometry. There is no spacetime interval and no light cone. The
degeneracy of $h$ ($\partial_t$ is in its kernel) is the geometric statement
that no spatial distance is defined between events at different times.

\subsection{Why the vector must be null}

If instead $\xi$ were spacelike or timelike, $g(\xi,\xi)\neq0$, the tensor
\begin{equation}
  \hat g=g-\frac{\xi^\flat\otimes\xi^\flat}{g(\xi,\xi)}
  \label{nr:KKtensor}
\end{equation}
would be basic and non-degenerate, and the reduction would be the familiar
Kaluza--Klein one: a Lorentzian (or Riemannian) metric on the quotient plus a
gauge field $\xi^\flat/g(\xi,\xi)$ (the connection form of the fibration). For a
null $\xi$ the denominator in \eqref{nr:KKtensor} vanishes and the construction
fails. What replaces the non-degenerate metric is the pair $(\tau,h)$ of
Proposition~\ref{nr:prop-induced}: the one-form $\xi^\flat$, which in
Kaluza--Klein theory is merely the gauge potential, becomes an independent
geometric datum (the clock), and the metric becomes degenerate. Nullness of
$\xi$ is what turns ``spacetime with a metric'' into ``absolute time with
Euclidean slices''.

\begin{nrdef}[Bargmann structure]
A \emph{Bargmann structure} is a triple $(\tilde M,g,\xi)$ with $g$ a Lorentzian
metric and $\xi$ a nowhere-vanishing null vector field with $\nabla\xi=0$
(Duval, Burdet, K\"unzle and Perrin 1985). Its quotient by the flow of $\xi$ is
a Newton--Cartan spacetime. What has been described above is the flat case.
\end{nrdef}

%---------------------------------------------------------------------
\section[Null geodesics and the origin of mass]{Dynamics: null geodesics and the origin of mass}
\label{sec:nr-dynamics}
%---------------------------------------------------------------------

\subsection{The Hamiltonian of null geodesics}

The phase space of a geodesic in $\tilde M$ is the cotangent bundle
$T^*\tilde M$, with canonical coordinates $(\mathbf{x},t,s)$ and conjugate
momenta $(\mathbf{p},p_t,p_s)$. Geodesics are the flows of the Hamiltonian
$\tilde{\mathcal{H}}=\tfrac12 g^{\mu\nu}p_\mu p_\nu$, and null geodesics lie
on the constraint surface $\tilde{\mathcal{H}}=0$. From \eqref{nr:metric},
\begin{equation}
  \tilde{\mathcal{H}}=\tfrac12\,\mathbf{p}^2+p_t\,p_s .
  \label{nr:Htilde}
\end{equation}
With an affine parameter $\lambda$, Hamilton's equations are
\begin{equation}
  \frac{\mathrm{d}\mathbf{x}}{\mathrm{d}\lambda}=\mathbf{p},\quad
  \frac{\mathrm{d}t}{\mathrm{d}\lambda}=p_s,\quad
  \frac{\mathrm{d}s}{\mathrm{d}\lambda}=p_t,\quad
  \frac{\mathrm{d}\mathbf{p}}{\mathrm{d}\lambda}=0,\quad
  \frac{\mathrm{d}p_t}{\mathrm{d}\lambda}=0,\quad
  \frac{\mathrm{d}p_s}{\mathrm{d}\lambda}=0 .
  \label{nr:hamiltonEq}
\end{equation}
The last equation is Noether's theorem for the Killing vector $\xi$: the
coordinate $s$ is cyclic, so its conjugate momentum $p_s$ is conserved. This
conserved quantity is the mass,
\begin{equation}
  m\;:=\;p_s .
  \label{nr:mass}
\end{equation}

\subsection{Recovering Newtonian motion}

Since $\mathrm{d}t/\mathrm{d}\lambda=m\neq0$, the coordinate $t$ is a good
parameter along the curve (take $m>0$ for future-directed motion). Define the
Galilean energy by $p_t=-E$. Then:
\begin{itemize}
\item The constraint $\tilde{\mathcal{H}}=0$ gives
      \begin{equation}
        E=\frac{\mathbf{p}^2}{2m},
        \label{nr:dispersion}
      \end{equation}
      the Newtonian free-particle dispersion relation.
\item $\mathrm{d}\mathbf{x}/\mathrm{d}t=\mathbf{p}/m$ and $\mathbf{p}=\text{const}$:
      uniform motion with $\mathbf{p}=m\mathbf{v}$.
\item The hidden coordinate evolves as
      \begin{equation}
        \frac{\mathrm{d}s}{\mathrm{d}t}=\frac{p_t}{p_s}=-\frac{E}{m}=-\tfrac12\mathbf{v}^2 .
        \label{nr:sdot}
      \end{equation}
\end{itemize}
Thus the projection $\pi$ of every null geodesic with $p_s=m\neq0$ is a
Newtonian free trajectory. The mass is not a parameter put into the geometry
by hand: all particles move in the same geometry \eqref{nr:metric}, and
different masses correspond to different values of a conserved momentum, much
as different charges in Kaluza--Klein theory correspond to different momenta
along the circle. A null geodesic with $p_s=0$ has $\mathrm{d}t/\mathrm{d}\lambda=0$
and, by the constraint, $\mathbf{p}=0$: it just slides along $\xi$ and does not
describe a Galilean particle.

There is one more reading of $m$. From \eqref{nr:hamiltonEq},
$m=\mathrm{d}t/\mathrm{d}\lambda$: the mass is the rate at which the absolute
time elapses per unit of the affine parameter of the null geodesic.

\subsection{The Newtonian action is $-m\,\Delta s$}

Integrating \eqref{nr:sdot} along a trajectory from $(\mathbf{x}_1,t_1)$ to
$(\mathbf{x}_2,t_2)$,
\begin{equation}
  \Delta s=-\int_{t_1}^{t_2}\tfrac12\mathbf{v}^2\,\mathrm{d}t
  =-\frac1m\int_{t_1}^{t_2}\tfrac12 m\mathbf{v}^2\,\mathrm{d}t
  =-\frac{S}{m},
  \qquad\text{i.e.}\qquad S=-m\,\Delta s .
  \label{nr:action}
\end{equation}
The hidden coordinate is minus the Newtonian action per unit mass, accumulated
along the motion. This one equation explains the quasi-invariance \eqref{nr:quasi}
geometrically. The ambient isometries (Section~\ref{sec:nr-symmetry}) act on the
\emph{whole} point $(\mathbf{x},t,s)$ and preserve the null curve, so they
preserve $\Delta s$ between two points of $\tilde M$. But the action of a boost
on the coordinate $s$ shifts it by a function $f(\mathbf{x},t)$ of the base
point, so when the two end points are expressed in terms of the base coordinates
alone, the Newtonian action $-m\Delta s$ changes by $-m[f(2)-f(1)]$: a boundary
term. This is the total derivative $\mathrm{d}F_{\mathbf u}/\mathrm{d}t$ of
\eqref{nr:quasi}, recovered as a coordinate effect in the fibre.

\subsection{Counting: the reduced phase space}

The phase space $T^*\tilde M$ has dimension $2(d+2)=2d+4$. The null-reduction
imposes two constraints that Poisson-commute with each other:
$p_s=m$ (generating translations of $s$) and $\tilde{\mathcal{H}}=0$
(generating reparametrisations along the geodesic). Both are first class, so
each removes two dimensions: one by the constraint itself, one by quotienting
along its flow. The reduced space has dimension $2d+4-4=2d$. Its points are
complete null geodesics modulo shifts of $s$, i.e.\ Newtonian motions, labelled
by $(\mathbf{x},\mathbf{p})$ at a fixed time. This is the usual phase space of a
Newtonian particle, obtained here as the \emph{space of motions} in the sense of
Souriau.

%---------------------------------------------------------------------
\section[The Bargmann group in the Poincar\'e group]{The Bargmann group inside the Poincar\'e group}
\label{sec:nr-symmetry}
%---------------------------------------------------------------------

\subsection{Killing vectors that commute with $\xi$}

The Killing vectors of the flat metric \eqref{nr:metric} are the generators
of the ambient Poincar\'e group $ISO(d+1,1)$,
\begin{equation}
  Y=a^\mu\partial_\mu+\omega^\mu{}_\nu\,x^\nu\partial_\mu,\qquad
  \omega_{\mu\nu}=g_{\mu\alpha}\omega^\alpha{}_\nu=-\omega_{\nu\mu},
  \label{nr:killing}
\end{equation}
with $x^\mu=(x^i,t,s)$ and $d+2+\tfrac12(d+2)(d+1)$ parameters. We ask for those
that preserve the null direction, $[\xi,Y]=0$. Since $\xi=\partial_s$ this means
that the components of $Y$ do not depend on $s$, i.e.\ $\omega^\mu{}_s=0$ for all
$\mu$, which in terms of the lowered components says
\begin{equation}
  \omega_{\mu s}=0\quad\text{for all }\mu
  \quad(\text{hence also }\omega_{s\mu}=0).
  \label{nr:omegacond}
\end{equation}
The surviving antisymmetric components are $\omega_{ij}$ (rotations, $\tfrac12 d(d-1)$
of them), $\omega_{it}=-\omega_{ti}$ ($d$ of them), and $\omega_{tt}=0$ by
antisymmetry. Raising an index with $g^{ts}=1$ gives $\omega^s{}_\nu=\omega_{t\nu}$
and $\omega^t{}_\nu=\omega_{s\nu}=0$: time is not rotated into anything, as
appropriate for absolute time. Writing $v_i=\omega_{it}$, the surviving vector
fields are
\begin{equation}
  P_i=\partial_i,\quad H=\partial_t,\quad M=\partial_s,\quad
  J_{ij}=x_i\partial_j-x_j\partial_i,\quad
  B_i=t\,\partial_i-x_i\,\partial_s .
  \label{nr:vf}
\end{equation}
Their number is
\begin{equation}
  \underbrace{d}_{P_i}+\underbrace{1}_{H}+\underbrace{1}_{M}
  +\underbrace{\tfrac12 d(d-1)}_{J_{ij}}+\underbrace{d}_{B_i}
  =\tfrac12 d(d-1)+2d+2 ,
  \label{nr:dim}
\end{equation}
which is $11$ for $d=3$. As a cross-check, the orbit of a null vector under the
Poincar\'e group of $\tilde M$ is the $(d+1)$-dimensional null cone, so the
stabiliser of $\xi$ has dimension
$\dim ISO(d+1,1)-(d+1)=\tfrac12(d+2)(d+3)-(d+1)$, which equals \eqref{nr:dim}:
for $d=3$, $15-4=11$.

\subsection{The algebra}

The non-vanishing commutators of the vector fields \eqref{nr:vf} are
\begin{equation}
  [B_i,P_j]=\delta_{ij}\,M,\qquad [B_i,H]=-P_i,\qquad
  [B_i,B_j]=0,\qquad [M,\,\cdot\,]=0,
  \label{nr:vfalgebra}
\end{equation}
together with the standard rotation relations. For example,
\[
  [B_i,P_j]=[t\partial_i-x_i\partial_s,\partial_j]=\delta_{ij}\,\partial_s .
\]
This is the Bargmann algebra \eqref{nr:KP}--\eqref{nr:KK}, but now every
generator is a genuine vector field on $\tilde M$ and the bracket is the
ordinary Lie bracket. The central element $M=\partial_s$ is central
\emph{by construction}: we imposed $[\xi,Y]=0$. In the language of the
introduction, the group of symmetries of the ambient space that preserve the
null vector is a central extension
\begin{equation}
  1\;\longrightarrow\;\mathbb{R}\xrightarrow{\;\;s\text{-shifts}\;\;}
  \text{Bargmann}\;\longrightarrow\;\text{Galilei}\;\longrightarrow\;1 ,
  \label{nr:exact}
\end{equation}
because the quotient by the $s$-shifts acts on $M$ as the Galilei group.

The conserved charges of null geodesics are $Q_Y=p_\mu Y^\mu$, and a short
calculation gives $\{Q_Y,Q_Z\}=-Q_{[Y,Z]}$. The sign is absorbed by defining
\begin{align*}
  P_i&=p_i=Q_{\partial_i}, & H&=-p_t=-Q_{\partial_t},\\
  M&=p_s=Q_{\partial_s}, & K_i&=x_ip_s-t\,p_i=-Q_{B_i},
\end{align*}
for which one finds directly the brackets \eqref{nr:KP}--\eqref{nr:KK} with
$m$ replaced by $M=p_s$, for instance
\[
  \{K_i,P_j\}=\{x_ip_s,p_j\}=p_s\,\delta_{ij}=M\,\delta_{ij}.
\]
On the surface $p_s=m$ these are exactly the Galilean charges of
Section~\ref{sec:nr-galilei}: $K_i=m\,x_i-t\,p_i$. The cocycle of
Section~\ref{sec:nr-galilei}, a constant $m\,\delta_{ij}$ that appeared as an
unexplained obstruction, is now the Poisson bracket of the two functions
$x_ip_s$ and $p_j$. It is an ordinary bracket in a bigger phase space.

\subsection{Finite transformations and the group commutator}

Exponentiating \eqref{nr:vf}, the Bargmann group acts on $\tilde M$ by
\begin{equation}
  \begin{aligned}
    \mathbf{x}&\mapsto R\,\mathbf{x}+\mathbf{v}\,t+\mathbf{a},\\
    t&\mapsto t+b,\\
    s&\mapsto s-\mathbf{v}\cdot R\,\mathbf{x}-\tfrac12 v^2\,t+c,
  \end{aligned}
  \label{nr:finite}
\end{equation}
with $R\in SO(d)$. That this preserves \eqref{nr:metric} follows from
$\mathrm{d}\mathbf{x}'=R\,\mathrm{d}\mathbf{x}+\mathbf{v}\,\mathrm{d}t$ and
$\mathrm{d}s'=\mathrm{d}s-\mathbf{v}\cdot R\,\mathrm{d}\mathbf{x}-\tfrac12 v^2\mathrm{d}t$:
\begin{align*}
  \mathrm{d}\mathbf{x}'^2+2\,\mathrm{d}t\,\mathrm{d}s'
  &=\bigl(\mathrm{d}\mathbf{x}^2+2\,\mathbf{v}\cdot R\,\mathrm{d}\mathbf{x}\,\mathrm{d}t+v^2\mathrm{d}t^2\bigr)\\
  &\quad+\bigl(2\,\mathrm{d}t\,\mathrm{d}s-2\,\mathbf{v}\cdot R\,\mathrm{d}\mathbf{x}\,\mathrm{d}t-v^2\mathrm{d}t^2\bigr)\\
  &=\mathrm{d}\mathbf{x}^2+2\,\mathrm{d}t\,\mathrm{d}s .
\end{align*}
Now compare a boost and a translation. Let $T_{\mathbf a}$ be $\mathbf{x}\mapsto
\mathbf{x}+\mathbf{a}$ and $B_{\mathbf v}$ the pure boost. Applying $T_{\mathbf a}$
first and then $B_{\mathbf v}$ gives
$(\mathbf{x},t,s)\mapsto(\mathbf{x}+\mathbf{a}+\mathbf{v}t,\;t,\;s-\mathbf{v}\cdot(\mathbf{x}+\mathbf{a})-\tfrac12 v^2t)$,
whereas the opposite order gives
$(\mathbf{x}+\mathbf{v}t+\mathbf{a},\;t,\;s-\mathbf{v}\cdot\mathbf{x}-\tfrac12 v^2t)$.
The two results agree in $(\mathbf{x},t)$ and differ in $s$ by $-\mathbf{v}\cdot\mathbf{a}$:
\begin{equation}
  B_{\mathbf v}\,T_{\mathbf a}\;=\;Z\bigl(-\mathbf{v}\cdot\mathbf{a}\bigr)\;T_{\mathbf a}\,B_{\mathbf v},
  \label{nr:groupcomm}
\end{equation}
where $Z(c)$ is the translation by $c$ in $s$, a central element. On the base
space, boosts and translations commute; on the ambient space their commutator is
a shift of the hidden coordinate by $-\mathbf{v}\cdot\mathbf{a}$. This is the
group-level form of $\{K_i,P_j\}=m\delta_{ij}$, and, acting on $e^{ims}$, it is
the phase $e^{-im\,\mathbf{v}\cdot\mathbf{a}}$ of Bargmann's superselection
argument.

\subsection{The Casimirs, seen from the ambient space}

The Poincar\'e Casimir of the ambient space,
$g^{\mu\nu}p_\mu p_\nu=\mathbf{p}^2+2p_tp_s=\mathbf{P}^2-2MH$ in terms of the
charges above, is invariant under the whole Poincar\'e group and in particular
under the Bargmann subgroup. Comparing with \eqref{nr:Casimirs},
\begin{equation}
  C\;=\;2MH-\mathbf{P}^2\;=\;-\,g^{\mu\nu}p_\mu p_\nu .
  \label{nr:CasimirAmbient}
\end{equation}
The Galilean internal-energy Casimir is minus the squared length of the
ambient momentum. For null geodesics $C=0$. A particle of ambient rest mass
$\mu$ has $C=\mu^2$, i.e.\ internal energy $U=\mu^2/2M$. We return to this in
Section~\ref{sec:nr-lightfront}.

%---------------------------------------------------------------------
\section{Newtonian gravity: the Eisenhart--Duval lift}
\label{sec:nr-gravity}
%---------------------------------------------------------------------

The flat Bargmann structure becomes a theory of dynamics with forces by
deforming the metric while keeping $\xi=\partial_s$. Take a function
$\Phi(\mathbf{x},t)$ and set
\begin{equation}
  g_\Phi=\delta_{ij}\,\mathrm{d}x^i\mathrm{d}x^j+2\,\mathrm{d}t\,\mathrm{d}s
        -2\,\Phi(\mathbf{x},t)\,\mathrm{d}t^2 .
  \label{nr:gPhi}
\end{equation}
Still $g(\xi,\xi)=0$ and nothing depends on $s$, so $\xi$ is null and Killing. Its
dual is still $\xi^\flat=\mathrm{d}t$, which is closed; a Killing vector whose
dual is closed is parallel, so $\nabla\xi=0$ and \eqref{nr:gPhi} is a
Bargmann structure. The inverse metric is
\begin{equation}
  g^{ij}=\delta^{ij},\qquad g^{ts}=1,\qquad g^{ss}=2\Phi,\qquad
  \text{all others }0 .
  \label{nr:ginv}
\end{equation}
(This is a plane-fronted wave with parallel rays, a Brinkmann pp-wave, and
this way of encoding Newtonian trajectories as geodesics of a metric one
dimension up is due to Eisenhart (1929); its identification with the Bargmann
structure is due to Duval, Burdet, K\"unzle and Perrin (1985) and
Duval, Gibbons and Horv\'athy (1991).)

\subsection{Null geodesics are Newtonian trajectories}

The null-geodesic Hamiltonian is
\begin{equation}
  \tilde{\mathcal{H}}=\tfrac12\,\mathbf{p}^2+p_tp_s+\Phi\,p_s^2 ,
  \label{nr:HtildePhi}
\end{equation}
and Hamilton's equations with affine parameter $\lambda$ read
\begin{equation}
  \begin{aligned}
  \frac{\mathrm{d}\mathbf{x}}{\mathrm{d}\lambda}&=\mathbf{p}, &
  \frac{\mathrm{d}t}{\mathrm{d}\lambda}&=p_s, &
  \frac{\mathrm{d}s}{\mathrm{d}\lambda}&=p_t+2\Phi\,p_s,\\
  \frac{\mathrm{d}\mathbf{p}}{\mathrm{d}\lambda}&=-p_s^2\,\nabla\Phi, &
  \frac{\mathrm{d}p_t}{\mathrm{d}\lambda}&=-p_s^2\,\partial_t\Phi, &
  \frac{\mathrm{d}p_s}{\mathrm{d}\lambda}&=0 .
  \end{aligned}
  \label{nr:hamiltonPhi}
\end{equation}
As before $p_s=m$ is conserved and $\mathrm{d}t/\mathrm{d}\lambda=m$. Dividing
by $m$ to pass to the parameter $t$, and setting $p_t=-E$:
\begin{align}
  \frac{\mathrm{d}\mathbf{x}}{\mathrm{d}t}&=\frac{\mathbf{p}}{m},
  &\frac{\mathrm{d}\mathbf{p}}{\mathrm{d}t}&=-m\,\nabla\Phi,
  &E&=\frac{\mathbf{p}^2}{2m}+m\,\Phi,
  \label{nr:newton}\\
  \frac{\mathrm{d}E}{\mathrm{d}t}&=m\,\partial_t\Phi,
  &\frac{\mathrm{d}s}{\mathrm{d}t}&=-\Bigl(\tfrac12\mathbf{v}^2-\Phi\Bigr)=-\frac{L}{m},
  &L&=\tfrac12 m\mathbf{v}^2-m\Phi .
  \label{nr:sL}
\end{align}
The first line is Newton's second law in the gravitational potential $\Phi$;
the second line says that along the motion $s$ is again minus the action per
unit mass, now for the full Lagrangian $L$. The relation $S=-m\Delta s$ of
\eqref{nr:action} holds with the interaction included. (The constraint
$\tilde{\mathcal{H}}=0$ gives the expression for $E$, and the expression for
$\mathrm{d}s/\mathrm{d}t$ uses $\mathrm{d}s/\mathrm{d}\lambda=p_t+2\Phi m$ with
$E=\mathbf{p}^2/2m+m\Phi$.)

\subsection{Universality of free fall}

The acceleration $\mathrm{d}^2\mathbf{x}/\mathrm{d}t^2=-\nabla\Phi$ in
\eqref{nr:newton} does not contain $m$. This is the equivalence principle, and
here it is automatic. All particles follow the null geodesics of one metric. A
null geodesic is fixed as a curve by a point and a direction, while its
affine parametrisation is free up to $\lambda\mapsto a\lambda+b$. The mass is
the value of $\mathrm{d}t/\mathrm{d}\lambda$, i.e.\ it merely records \emph{which}
affine scale is used. Two particles of different masses starting at the same
point with the same velocity trace out the same curve; they differ only in the
normalisation of the tangent vector, and hence of the momentum
$\mathbf{p}=m\mathbf{v}$. The statement ``gravitational and inertial mass are
equal'' becomes the statement that the same number $p_s$ is both the charge
that couples the particle to $g_\Phi$ and the factor relating its momentum to
its velocity.

\subsection{Field equation}

The Ricci tensor of $g_\Phi$ has a single non-zero component. A direct
computation gives
\begin{equation}
  R_{\mu\nu}=\bigl(\nabla^2\Phi\bigr)\,\tau_\mu\tau_\nu,\qquad \tau=\mathrm{d}t,
  \label{nr:Ricci}
\end{equation}
where $\nabla^2$ is the spatial Laplacian. With the convention in which
positive-energy matter satisfies $R_{\mu\nu}k^\mu k^\nu\ge0$ for null $k$,
Poisson's equation $\nabla^2\Phi=4\pi G\rho$ is therefore equivalent to
\begin{equation}
  R_{\mu\nu}=4\pi G\,\rho\;\tau_\mu\tau_\nu ,
  \label{nr:field}
\end{equation}
an Einstein-type equation in the ambient space. For curved
Newton--Cartan geometries with closed clock form, $\mathrm{d}\tau=0$, Duval et
al.\ showed that the theory is described by Bargmann structures subject to an
equation of this type. For geometries with gravitational time dilation
($\mathrm{d}\tau\neq0$), the Newton--Cartan equations coupled to matter are not
obtained from a null reduction in this simple way (see the 2023 review of
non-relativistic gravity, arXiv:2212.11309).

\begin{nrrem}[Gauge fields and non-inertial frames]
Adding a one-form $A=A_i\mathrm{d}x^i$ to the metric,
$g=\delta_{ij}\mathrm{d}x^i\mathrm{d}x^j+2\,\mathrm{d}t(\mathrm{d}s+A_i\mathrm{d}x^i)-2U\,\mathrm{d}t^2$,
leads to $g^{si}=-A_i$, $g^{ss}=2U+\mathbf{A}^2$ and to the
Hamiltonian $\tfrac12(\mathbf{p}-p_s\mathbf{A})^2+p_tp_s+Up_s^2$, i.e.\ to
$E=(\mathbf{p}-m\mathbf{A})^2/2m+mU$. This is minimal coupling with a coupling
constant proportional to $m$: the same universality of the
charge-to-mass ratio, which is why a term of this kind describes Coriolis
and magnetic-type forces that are felt equally by every mass, but not
a general electromagnetic coupling with an independent charge.
\end{nrrem}

%---------------------------------------------------------------------
\section{The quantum epilogue}
\label{sec:nr-quantum}
%---------------------------------------------------------------------

Set $\hbar=1$. The massless wave equation on $\tilde M$ is
\begin{equation}
  g^{\mu\nu}\partial_\mu\partial_\nu\Phi=\bigl(\nabla^2+2\,\partial_t\partial_s\bigr)\Phi=0 .
  \label{nr:wave}
\end{equation}
Because $\partial_s$ is a symmetry, one may Fourier-decompose
\begin{equation}
  \Phi=\int\mathrm{d}m\;e^{\,ims}\,\psi_m(\mathbf{x},t),
\end{equation}
and each sector evolves independently. For a single sector, $\partial_s\to im$ and \eqref{nr:wave} becomes
\begin{equation}
  i\,\partial_t\psi=-\frac{1}{2m}\nabla^2\psi ,
  \label{nr:schrodinger}
\end{equation}
the free Schr\"odinger equation, with $m$ the eigenvalue of $-i\partial_s$. If
the wave equation carries a mass term, $(g^{\mu\nu}\partial_\mu\partial_\nu-\mu^2)\Phi=0$,
the same reduction gives $i\partial_t\psi=(-\nabla^2+\mu^2)\psi/2m$, i.e.\ an
internal energy $U=\mu^2/2m$, in agreement with \eqref{nr:CasimirAmbient}.

\subsection{Where the boost phase comes from}

The ambient field transforms as a scalar under the isometry \eqref{nr:finite},
$\Phi'(x')=\Phi(x)$. Take a pure boost, $\mathbf{x}'=\mathbf{x}+\mathbf{v}t$,
$s'=s-\mathbf{v}\cdot\mathbf{x}-\tfrac12v^2t$. Then
$\psi'(\mathbf{x}',t)\,e^{ims'}=\psi(\mathbf{x},t)\,e^{ims}$ gives
\begin{equation}
  \psi'(\mathbf{x}',t)
  =e^{\,im(\mathbf{v}\cdot\mathbf{x}+\frac12v^2t)}\psi(\mathbf{x},t)
  =e^{\,im(\mathbf{v}\cdot\mathbf{x}'-\frac12v^2t)}\,\psi(\mathbf{x}'-\mathbf{v}t,\,t).
  \label{nr:boostphase}
\end{equation}
This is the familiar Galilean transformation of the wave function. In the usual
presentation the phase $e^{im(\mathbf{v}\cdot\mathbf{x}-\frac12v^2t)}$ is
added by hand to make the Schr\"odinger equation covariant, and one says that
the Galilei group acts only projectively. Here the phase is the plain geometric
statement that the boost shifts the coordinate $s$ by
$-\mathbf{v}\cdot\mathbf{x}-\tfrac12v^2t$, which is the same function
$f(\mathbf{x},t)$ that produced the total derivative in \eqref{nr:quasi}. The
group commutator \eqref{nr:groupcomm} acts on $e^{ims}$ as multiplication by
$e^{-im\,\mathbf{v}\cdot\mathbf{a}}$, which is the phase of Bargmann's
superselection argument in Section~\ref{sec:nr-galilei}.

\begin{nrrem}[Quantised mass]
If the null coordinate is made periodic, $s\sim s+2\pi R$ (a null
identification, the starting point of discrete light-cone quantisation), then
$e^{ims}$ is single valued only for $m=n/R$, $n\in\mathbb{Z}$. Mass becomes
quantised in integer units of $1/R$, and additivity of the mass becomes
additivity of the integer $n$ (``number of quanta''). This is the logic by
which particle number appears in discrete light-cone quantisation.
\end{nrrem}

%---------------------------------------------------------------------
\section[Relativistic light-front dynamics]{The same mathematics inside relativistic dynamics: the light front}
\label{sec:nr-lightfront}
%---------------------------------------------------------------------

The construction does not require inventing a fifth dimension. Take $d=2$ in
\eqref{nr:metric}. Then $\tilde M=\mathbb{R}^{2,1}\times\mathbb{R}$ is the
ordinary four-dimensional Minkowski space $\mathbb{R}^{3,1}$ with coordinates
$(x^1,x^2,t,s)$ related to the Minkowski coordinates by \eqref{nr:nullcoords}
(with $X^3\equiv X^{d+1}$ and $d+1=3$): in the usual terminology of light-front
quantisation,
\[
  t=x^+=\frac{X^0+X^3}{\sqrt2},\qquad s=-x^-=\frac{X^3-X^0}{\sqrt2},
\]
and the transverse plane is $(x^1,x^2)=\mathbf{x}_\perp$. A direct computation
of the conjugate momenta gives
\begin{equation}
  M=p_s=P^+=\frac{P^0+P^3}{\sqrt2},\qquad
  H=-p_t=P^-=\frac{P^0-P^3}{\sqrt2}.
  \label{nr:lfmomenta}
\end{equation}
Thus the \emph{light-front momentum $P^+$ plays the role of the Galilean mass}
and the \emph{light-front Hamiltonian $P^-$ plays the role of the Galilean
energy}, in a $2{+}1$-dimensional Galilean world with time $x^+$.

\subsection{Relativistic mass is the Bargmann Casimir}

A relativistic particle of rest mass $\mu$ in $\mathbb{R}^{3,1}$ satisfies
$g^{\mu\nu}p_\mu p_\nu=-\mu^2$, that is $\mathbf{p}_\perp^2+2p_tp_s=-\mu^2$, or
\begin{equation}
  P^-=\frac{\mathbf{p}_\perp^2+\mu^2}{2P^+}.
  \label{nr:lfdispersion}
\end{equation}
This is a Galilean particle of mass $M=P^+$ in two space dimensions whose
internal energy is $U=\mu^2/2M$. Comparing with \eqref{nr:CasimirAmbient},
\begin{equation}
  \underbrace{2MH-\mathbf{P}^2}_{\text{Bargmann Casimir } C}
  \;=\;\underbrace{\mu^2}_{\text{Poincar\'e Casimir}} .
  \label{nr:twoCasimirs}
\end{equation}
The two roles of mass of the introduction meet in a single equation. The
Poincar\'e Casimir $\mu^2$ of the $3{+}1$-dimensional theory is, from the point
of view of the transverse $2{+}1$ Galilean subtheory, the combination
$2MH-\mathbf{P}^2$ (twice the mass times the internal energy), while the
\emph{central charge} of the Galilean subalgebra is the light-front momentum
$P^+$.

\subsection{The Bargmann group inside the Poincar\'e group}

Dirac (1949) classified the ``forms of relativistic dynamics'' by the subgroup
of the Poincar\'e group that leaves a hypersurface invariant. For the light
front $x^+=0$ this \emph{kinematical subgroup} has seven generators in
$3{+}1$ dimensions: $P^+$, $P_1$, $P_2$, $J_{12}$, the two transverse boosts
$E_i=J_{+i}$, and the longitudinal boost $K_3$; the Hamiltonian $P^-$ is
\emph{dynamical}. The Bargmann group of $d=2$ (which has $\tfrac12\cdot2\cdot1+4+2=7$
generators by \eqref{nr:dim}) consists of $P^+$, $P_1$, $P_2$, $J_{12}$, $E_1$, $E_2$
and $P^-$. It differs from the front-form kinematical subgroup only in that
$P^-$ is included and the longitudinal boost $K_3$ is not: $K_3$ rescales
$\xi$, $(t,s)\mapsto(\lambda t,s/\lambda)$, so it does not stabilise the vector
$\xi$ (it only preserves its direction), and it acts by $(M,H)\mapsto(\lambda M,H/\lambda)$.
The transverse boosts $E_i$ satisfy $\{E_i,P_j\}=\delta_{ij}P^+$ and are
\emph{exactly} Galilean boosts. This is why light-front relativistic dynamics
has an exact Galilean invariance in the transverse plane, and why the
``infinite momentum frame'' calculations of the 1960s look nonrelativistic.

\begin{nrrem}[Which spacetime is the parent?]
The null reduction relates a Galilean world in $d$ space dimensions to a
relativistic world in $d+1$ space dimensions, not to the relativistic world in
$d$ space dimensions. The physical $3{+}1$ Galilean particle is a sector of
massless motion in $4{+}1$ Minkowski space; the light-front picture of an
actual $3{+}1$ relativistic particle gives a $2{+}1$ Galilean particle. We keep
this distinction in view in the next section.
\end{nrrem}

%---------------------------------------------------------------------
\section{Extension, contraction, reduction}
\label{sec:nr-limit}
%---------------------------------------------------------------------

There are now three distinct ways in which the Bargmann algebra, and with it the
Galilean mass, is connected to the rest of mechanics. They are often run
together, so we separate them.

\subsection{The contraction route}

Restore the speed of light. The relativistic free particle of rest mass $\mu$ has
$E=\sqrt{\mathbf{p}^2c^2+\mu^2c^4}$ and boost charge
$\mathbf{K}=E\,\mathbf{x}/c^2-t\,\mathbf{p}$. A direct calculation gives the
Poisson algebra
\begin{equation}
  \{K_i,P_j\}=\delta_{ij}\,\frac{E}{c^2},\qquad
  \{K_i,E\}=P_i,\qquad
  \{K_i,K_j\}=-\frac{1}{c^2}\,\epsilon_{ijk}J_k ,
  \label{nr:poincare}
\end{equation}
with $\{P_i,E\}=0$ and the usual rotation relations. If one simply lets
$c\to\infty$ in the structure constants, the right-hand sides of the first and
third relations vanish, and one obtains the Galilei algebra \eqref{nr:galilei}
\emph{without} a central term. The mass has disappeared, because the
bracket $\{K_i,P_j\}$ is $E/c^2$, and $E\to\infty$ as $c\to\infty$.

To keep the central term one must keep track of the large part of the
energy. Adjoin to the Poincar\'e algebra a generator $Z$ that commutes with
everything (this is an extra input, it is not present in the Poincar\'e algebra)
and set
\begin{equation}
  E=c^2Z+H' .
  \label{nr:Esplit}
\end{equation}
The relations \eqref{nr:poincare} become
\begin{equation}
  \{K_i,P_j\}=\delta_{ij}\Bigl(Z+\frac{H'}{c^2}\Bigr),\qquad
  \{K_i,H'\}=P_i,\qquad
  \{K_i,K_j\}=-\frac{\epsilon_{ijk}J_k}{c^2}.
  \label{nr:contraction}
\end{equation}
In the limit $c\to\infty$ at fixed $Z$, $H'$, $\mathbf{P}$, $\mathbf{K}$, $\mathbf{J}$ these become
\begin{equation}
  \{K_i,P_j\}=Z\,\delta_{ij},\qquad \{K_i,H'\}=P_i,\qquad \{K_i,K_j\}=0 ,
  \label{nr:contractionlimit}
\end{equation}
which is the Bargmann algebra with $M=Z$. On an irreducible representation
of mass $\mu$ one takes $Z=\mu$; then
\begin{equation}
  \begin{aligned}
  H'&=E-\mu c^2=\frac{\mathbf{p}^2}{2\mu}-\frac{\mathbf{p}^4}{8\mu^3c^2}+\cdots
     \;\longrightarrow\;\frac{\mathbf{p}^2}{2\mu},\\
  \mathbf{K}&=\Bigl(\mu+\frac{H'}{c^2}\Bigr)\mathbf{x}-t\,\mathbf{p}
     \;\longrightarrow\;\mu\,\mathbf{x}-t\,\mathbf{p}.
  \end{aligned}
  \label{nr:expansion}
\end{equation}
This is the sense of the statement ``the central charge is the large-$c$ limit
of the Casimir (rest energy)$/c^2$'': it is exact on each irreducible
representation, but at the level of algebras it requires the extra generator
$Z$. Moreover the limit is singular in the technical sense of a
contraction (In\"on\"u and Wigner 1953; Bacry and L\'evy-Leblond 1968): for every
finite $c$ the algebra is Poincar\'e$\oplus\mathbb{R}$, containing the
simple Lorentz algebra $\mathfrak{so}(3,1)$, while the limit contains the
abelian boosts, an algebra of a different isomorphism class. The relativistic
algebra is a deformation of the Bargmann algebra with parameter $1/c^2$, not
the other way round.

\subsection{Three routes compared}

\begin{table}[ht]
\centering
\footnotesize
\setlength{\tabcolsep}{4pt}
\renewcommand{\arraystretch}{1.25}
\begin{tabular}{@{}p{0.13\linewidth}p{0.29\linewidth}p{0.20\linewidth}p{0.27\linewidth}@{}}
\toprule
\textbf{Route} & \textbf{Operation} & \textbf{Mass appears as} & \textbf{Relation between the two kinematics}\\
\midrule
Extension &
Central extension of the Galilei algebra by the cocycle $\Theta(K_i,P_j)=m\delta_{ij}$ &
Central charge $M$, fixed by hand in the extension &
Internal to Galilean mechanics; no relativity involved\\
Contraction &
$c\to\infty$ of the Poincar\'e algebra plus an adjoined central generator, $E=c^2Z+H'$ &
Casimir $\mu$, which becomes the value of the central charge $Z$ &
Limit of a family of algebras (singular, expansion in $1/c^2$); same spacetime dimension\\
Reduction &
Quotient of $\mathbb{R}^{d+1,1}$ by a null parallel vector $\xi$; Bargmann group $=$ stabiliser of $\xi$ &
Momentum $p_s$ conjugate to the null coordinate &
Exact, no parameter and no error terms; the parent spacetime has one more space dimension\\
\bottomrule
\end{tabular}
\caption{Three connections between the Bargmann algebra and other structures in
mechanics.}
\label{tab:nr-routes}
\end{table}

The reduction differs in kind from the contraction in three respects.
First, it produces the central generator instead of requiring it as an input:
$M=\partial_s$ exists because the ambient space has a coordinate $s$.
Second, it contains no limit and no small parameter: the Newtonian
dispersion $E=\mathbf{p}^2/2m$ in \eqref{nr:dispersion} is exact for null
geodesics, whereas \eqref{nr:expansion} is an asymptotic expansion. Third, the
relativistic parent is a different spacetime. The contraction compares the
Galilean theory in $d$ space dimensions with the relativistic theory in the
\emph{same} $d$ dimensions. The reduction compares it with the
\emph{massless} sector of a relativistic theory in $d+1$ dimensions, and
Section~\ref{sec:nr-lightfront} shows that a massive relativistic particle in
$d+1$ dimensions appears as a Galilean particle in $d$ dimensions with internal
energy $\mu^2/2M$. The two relations are compatible, since both lead to the
same Bargmann algebra, but they say different things.

%---------------------------------------------------------------------
\section[Klein--Rumer five-dimensional optics]{The spacelike precursor: Klein and Rumer's five-dimensional optics}
\label{sec:nr-rumer}
%---------------------------------------------------------------------

The idea that mechanics is optics in one more dimension, with the extra
coordinate related to the action, is older than the Bargmann construction. It
goes back to Hamilton's optical--mechanical analogy and, in its five-dimensional
form, to Klein (1926), who took a \emph{massless} wave in five dimensions and
obtained the Hamilton--Jacobi equation of a four-dimensional particle in the ray
limit. Rumer developed it as ``5-optics'' (1949, 1956): the action is
treated as a fifth coordinate, so that a mechanics problem becomes a problem of
ray propagation in a space of coordinates, time and action, with the geometrical
5-optics the ray limit of a wave 5-optics. The programme was continued, in a
different mathematical setting, by Keller's START (space--time--action) theory.
In this section we compare it with the null reduction. We use Klein's form of
the construction (units $c=1$); the details of Rumer's own formulation, which we
know only from secondary sources, are not reproduced, and the reader should
consult the primary works listed at the end of the chapter.

%.....................................................................
\subsection{Eikonal in an extended space}
%.....................................................................

Let $\tilde M$ carry a metric $g$ and a Killing vector $\xi$ with coordinate $y$
along it, so $\xi=\partial_y$. A wave $\Psi=e^{i\Theta/\hbar}$ satisfying
$\Box\Psi=0$ has, in the ray limit $\hbar\to0$, the eikonal equation
$g^{AB}\partial_A\Theta\,\partial_B\Theta=0$: surfaces of constant phase are
null hypersurfaces and rays are null geodesics. If $\Theta=p_y\,y+S(x^\mu)$,
with $p_y$ the (conserved) momentum along $\xi$, the eikonal equation becomes a
Hamilton--Jacobi equation for $S$ on the quotient, with $p_y$ appearing as a
parameter. The two cases to compare differ only in the nature of $\xi$.

\paragraph{Case A: spacelike $\xi$ (Klein, Rumer).}
Take $g=-\mathrm{d}t^2+\mathrm{d}\mathbf{x}^2+\mathrm{d}y^2$ and
$\Theta=m\,y+S(\mathbf{x},t)$ with $m=p_y>0$. The eikonal equation is
\begin{equation}
  -(\partial_tS)^2+(\nabla S)^2+m^2=0 ,
  \label{nr:HJrel}
\end{equation}
the relativistic Hamilton--Jacobi equation of a particle of mass $m$: the rest
mass is the momentum along the extra direction. On a null ray,
$\mathrm{d}y^2=\mathrm{d}t^2-\mathrm{d}\mathbf{x}^2=\mathrm{d}\tau^2$, so $y$
advances along the ray by the proper time, and the relativistic free action is
\begin{equation}
  S_{\rm rel}=-m\int\mathrm{d}\tau=-m\,\Delta y .
  \label{nr:Srel}
\end{equation}

\paragraph{Case B: null $\xi$ (this chapter).}
Take $g=\mathrm{d}\mathbf{x}^2+2\,\mathrm{d}t\,\mathrm{d}s$ and
$\Theta=m\,s+S(\mathbf{x},t)$. The eikonal equation is
\begin{equation}
  \partial_tS+\frac{(\nabla S)^2}{2m}=0 ,
  \label{nr:HJGal}
\end{equation}
the free Newtonian Hamilton--Jacobi equation (with $\Phi$ added as in
\eqref{nr:gPhi}, one gets $\partial_tS+(\nabla S)^2/2m+m\Phi=0$), and
$S=-m\,\Delta s$ by \eqref{nr:action}.

\begin{table}[ht]
\centering
\footnotesize
\setlength{\tabcolsep}{4pt}
\renewcommand{\arraystretch}{1.25}
\begin{tabular}{@{}p{0.20\linewidth}p{0.36\linewidth}p{0.36\linewidth}@{}}
\toprule
 & \textbf{Klein--Rumer (spacelike)} & \textbf{Bargmann (null)}\\
\midrule
Extra direction $\xi$ & $g(\xi,\xi)>0$ & $g(\xi,\xi)=0$\\
Left on the quotient & Lorentzian metric plus a gauge field & clock form and degenerate spatial metric\\
Mass & $p_y$, with $E^2=\mathbf{p}^2+p_y^2$ & $p_s$, with $E=\mathbf{p}^2/2p_s$\\
Hamilton--Jacobi equation & \eqref{nr:HJrel} & \eqref{nr:HJGal}\\
Action along a ray & $-m\,\Delta y$, $y=$ proper time & $-m\,\Delta s$\\
\bottomrule
\end{tabular}
\caption{The same eikonal construction with a spacelike and a null extra direction.}
\label{tab:nr-rumer}
\end{table}

%.....................................................................
\subsection{What the two constructions share}
%.....................................................................

\paragraph{The action is mass times a coordinate.}
In both cases the action is $-m$ times the displacement of the extra coordinate,
$S=-m\,\Delta\chi$ with $\chi=y$ or $\chi=s$. If the fifth coordinate
is taken to be the action itself, as in Rumer's 5-optics, then $x^5=S=-m\chi$,
and the coordinate of a particle depends on its mass. If one divides
by $m$, as here, the coordinate is the action per unit mass and \emph{all
particles live in one geometry}, with the mass appearing as a conserved
momentum.

\paragraph{The action is defined up to a function.}
Rumer emphasised that the action is defined only up to $S\mapsto S+f(x^\mu)$.
In the spacelike case, with $g=g_{\mu\nu}\mathrm{d}x^\mu\mathrm{d}x^\nu+(\mathrm{d}y+A_\mu\mathrm{d}x^\mu)^2$,
the redefinition $y\mapsto y+\lambda(x)$ is the gauge transformation
$A_\mu\mapsto A_\mu-\partial_\mu\lambda$. In the null case it is the same
freedom, $s\mapsto s+\lambda(\mathbf{x},t)$, with the same gauge-transformation
interpretation (Section~\ref{sec:nr-applications}); the Galilean boost is the
particular choice $\lambda=-\mathbf{v}\cdot\mathbf{x}-\tfrac12v^2t$, and the
cocycle of Section~\ref{sec:nr-galilei} is its trace.
In both constructions the field $A$ couples with strength proportional to the
momentum along $\xi$, so the charge-to-mass ratio is fixed by the geometry. This
is the well-known obstacle to reading $A_\mu$ in Kaluza--Klein theory as the
electromagnetic potential of ordinary particles with independent charges; in
the null case the same universality is the equivalence principle, and one reads
$A$ as a gravito-inertial potential (Section~\ref{sec:nr-gravity}).

\paragraph{Ray and wave.}
The geometrical (ray) limit of the five-dimensional massless wave equation is
\eqref{nr:HJrel} in Case~A and \eqref{nr:HJGal} in Case~B; the wave equation
itself gives the Klein--Gordon equation and the Schr\"odinger equation
(Section~\ref{sec:nr-quantum}), respectively. This is Rumer's ``geometrical 5-optics
as a limit of wave 5-optics'' in both settings.

%.....................................................................
\subsection{One light cone, two slicings}
%.....................................................................

Both cases are projections of massless motion in the same flat space
$\mathbb{R}^{d+1,1}$, differing in which Killing direction is used to form the
quotient. A spacelike quotient gives relativistic mechanics of massive particles
in $d$ space dimensions; a null quotient gives Newtonian mechanics. The two are
related at the level of the action. Along a Case~A ray, expanding
$\mathrm{d}\tau=\sqrt{1-v^2}\,\mathrm{d}t$,
\begin{equation}
  -m\,\Delta(y-t)=S_{\rm rel}+m\,\Delta t
  =\int\Bigl(\tfrac12mv^2+\tfrac18mv^4+\cdots\Bigr)\mathrm{d}t ,
  \label{nr:restsub}
\end{equation}
whose leading term is the Galilean action $-m\,\Delta s$ of Case~B. Replacing $y$
by $y-t$, which is the null coordinate up to normalisation, subtracts the rest
action $-m\,\Delta t$, in the same way as the substitution
$\Psi=e^{-imt}\psi$ takes the Klein--Gordon equation to the Schr\"odinger equation
and as the split $E=c^2Z+H'$ does in \eqref{nr:Esplit}. This is a correspondence at
leading order, not an identification: the exact null time is $(t+y)/\sqrt2$,
which differs from $t$ at the order of the relativistic corrections, and the
null reduction has no corrections because it is a different quotient, not an
expansion of the first. The precise sense in which a null direction is a limit
of boosted spacelike directions is the one used in discrete light-cone
quantisation (Seiberg 1997).

%.....................................................................
\subsection{Status and an objection}
%.....................................................................

Rumer's 5-optics remains a niche tradition. What survives in the mainstream is
its structure, under other names: Kaluza--Klein reduction (Case~A), and the
Eisenhart--Duval lift and Bargmann structures (Case~B). Keller's START theory
continues the idea with the five-dimensional homogeneous space of the conformal
group $O(4,2)$ (spaces of anti-de Sitter or de Sitter type, bounded by the
conformally compactified Minkowski space), a geometry that resembles the one used in
holography; we note the resemblance but do not claim any established
relation to Section~\ref{sec:nr-holography}.

One objection raised in the literature (Redington 1997) is that a quantity
identified with a coordinate should not behave as an extensive variable, which
would exclude the action. The extra coordinate $s$ used here is conjugate to the
extensive total mass, so it is intensive, in the way the chemical potential is
conjugate to the particle number; the action $S=-m\,s$ is extensive. This is an
observation about the present construction and is not taken from the cited work.
It also fits the reading of Section~\ref{sec:nr-observable}: $s$ is a bookkeeping
coordinate, and the physical content is the conserved extensive quantity $m$.

%---------------------------------------------------------------------
\section{What is observable in the construction}
\label{sec:nr-observable}
%---------------------------------------------------------------------

The reduction can be read in two ways, and it is worth stating which parts of
the geometry are accessible to measurement and which are not.

The base $M$ and its structure $(\tau,h)$ are what the Newtonian theory
speaks about: clocks, rulers, trajectories $\mathbf{x}(t)$, momenta
$\mathbf{p}=m\mathbf{v}$, and the conserved mass $m=p_s$, which is measured by
inertia, for example by the response to a standard force. All of these are
functions on the base or on the reduced phase space of Section~\ref{sec:nr-dynamics}.

The coordinate $s$ is not of this kind. By \eqref{nr:action} it is minus the
Newtonian action per unit mass, accumulated along a motion. The action is
defined only up to boundary terms and, as the boost transformation
\eqref{nr:finite} shows, even the \emph{value} of the difference $\Delta s$
between two points of the base depends on the frame in which the end
points are described, because a boost shifts $s$ by a function of position and
time. No instrument measures $s$. What can be measured is its conjugate momentum,
and the way in which $s$ enters physics is through the quasi-invariance of the
action and the phase $e^{ims}$ of quantum states. In this reading, the extra
dimension is an organising device: it packages the cocycle of
Section~\ref{sec:nr-galilei} into a geometry, and the physical content is
entirely in the quotient together with the conserved number $m$.

The relation between the Galilean theory and the ambient relativistic theory
is therefore of the type \emph{projection}: $\pi:\tilde M\to M$ forgets the
metric and the coordinate $s$, and keeps the clock, the spatial metric, and the
momentum along the forgotten direction. It is exact, in the sense that the
projected null geodesics are \emph{precisely} Newtonian trajectories, and it
differs in this respect from a limiting relation, which holds only up to terms
that vanish in the limit. It is a statement about the structure of the theory,
not a claim that the Newtonian particle ``really'' moves in five dimensions.

%---------------------------------------------------------------------
\section[Applications]{Applications: where the construction is used}
\label{sec:nr-applications}
%---------------------------------------------------------------------

Up to now the null reduction has reorganised facts about Galilean mechanics
that were already known. This section shows three places where the
reorganisation does real work in current research. They have a common
mechanism. A Galilean system coupled to background fields is the same thing as a
massless field on a Bargmann space; the extra coordinate then supplies
the particle-number symmetry and the Galilean boosts automatically, and the
background fields are components of one higher-dimensional metric. We work out
one example in each case. Results that are quoted from the literature
rather than derived here are marked as such.

%.....................................................................
\subsection{Coupling a Galilean field to background fields}
%.....................................................................

A Galilean-invariant quantum field theory is easy to write in flat space and
awkward to put on a curved or time-dependent background: the Lorentz-covariant
recipe of replacing $\eta_{\mu\nu}$ by $g_{\mu\nu}$ is not available, because
there is no spacetime metric (Section~\ref{sec:nr-geometry}). The null reduction
supplies the recipe. Take the Bargmann metric with the gauge field and the
potential of Section~\ref{sec:nr-gravity},
\begin{equation}
  g_{A,U}=\delta_{ij}\,\mathrm{d}x^i\mathrm{d}x^j
  +2\,\mathrm{d}t\,\bigl(\mathrm{d}s+A_i\,\mathrm{d}x^i\bigr)-2U\,\mathrm{d}t^2 ,
  \label{nr:gAU}
\end{equation}
with $A_i(\mathbf{x},t)$ and $U(\mathbf{x},t)$ arbitrary. Its determinant is $\det g=-1$, and its inverse
has the components
\begin{equation}
  g^{ij}=\delta^{ij},\quad g^{ts}=1,\quad g^{si}=-A_i,\quad
  g^{ss}=\mathbf{A}^2+2U,\quad\text{all others }0 .
  \label{nr:ginvAU}
\end{equation}

\begin{nrprop}\label{nr:prop-reducedaction}
Let $\Phi$ be a complex scalar on $(\tilde M,g_{A,U})$ with the massless action
$S[\Phi]=-\int\mathrm{d}^dx\,\mathrm{d}t\,\mathrm{d}s\,\sqrt{-g}\,g^{\mu\nu}\partial_\mu\bar\Phi\,\partial_\nu\Phi$.
For $\Phi=e^{\,ims}\psi(\mathbf{x},t)$ the action per unit length of $s$, divided by $2m$, is
\begin{equation}
  \begin{aligned}
  S[\psi]&=\int\mathrm{d}t\,\mathrm{d}^dx\;\Bigl[
  \frac{i}{2}\bigl(\bar\psi\,\partial_t\psi-\psi\,\partial_t\bar\psi\bigr)
  -\frac{1}{2m}\,\bigl|D_i\psi\bigr|^2-m\,U\,|\psi|^2\Bigr],\\
  D_i&=\partial_i-i\,m\,A_i .
  \end{aligned}
  \label{nr:schLag}
\end{equation}
\end{nrprop}
\begin{proof}
Since $\sqrt{-g}=1$, only $g^{\mu\nu}\partial_\mu\bar\Phi\partial_\nu\Phi$ matters. With $\partial_s\Phi=im\Phi$
and $\partial_s\bar\Phi=-im\bar\Phi$, the four types of term in \eqref{nr:ginvAU} give
\begin{align*}
  g^{ij}&:\quad \partial_i\bar\psi\,\partial_i\psi,\\
  g^{ts}&:\quad im\bigl(\psi\,\partial_t\bar\psi-\bar\psi\,\partial_t\psi\bigr),\\
  g^{si}&:\quad -im\,A_i\bigl(\psi\,\partial_i\bar\psi-\bar\psi\,\partial_i\psi\bigr),\\
  g^{ss}&:\quad m^2\bigl(\mathbf{A}^2+2U\bigr)|\psi|^2 .
\end{align*}
The first, third and the $\mathbf{A}^2$ part of the fourth combine into
$|D_i\psi|^2=|\partial_i\psi|^2+imA_i(\bar\psi\partial_i\psi-\psi\partial_i\bar\psi)+m^2\mathbf{A}^2|\psi|^2$.
Hence $-g^{\mu\nu}\partial_\mu\bar\Phi\partial_\nu\Phi=-|D_i\psi|^2+im(\bar\psi\partial_t\psi-\psi\partial_t\bar\psi)-2m^2U|\psi|^2$,
and dividing by $2m$ gives \eqref{nr:schLag}.
\end{proof}

Equation \eqref{nr:schLag} is the Schr\"odinger Lagrangian for a particle of
mass $m$ in a potential $mU$ and a gauge field $mA_i$: \emph{the mass is the
charge of the $U(1)$ that comes from translations of $s$}. Three consequences
follow at once.
\begin{enumerate}
\item \emph{Gauge invariance.} A change of fibre coordinate
$s\mapsto s'=s+\lambda(\mathbf{x},t)$ preserves the form of $g_{A,U}$ with
$A_i\mapsto A_i-\partial_i\lambda$ and $U\mapsto U+\partial_t\lambda$, and
$\Phi'(x')=\Phi(x)$ gives $\psi\mapsto e^{-im\lambda}\psi$. This is the usual
gauge transformation with charge $m$ and gauge potential $(A_i,-U)$.
\item \emph{Boosts are free.} For $A=U=0$ the metric is flat and the
Galilean boosts \eqref{nr:finite} are isometries; \eqref{nr:boostphase}
is then their action on $\psi$. In the curved case, any diffeomorphism of
$\tilde M$ that preserves $\xi$ acts as a symmetry, with no further input.
\item \emph{Universality.} The same $A_i$ and $U$ couple to every mass in
the proportion $m$, so these background fields describe gravitational and
inertial effects (potentials, Coriolis forces, accelerating frames), not
arbitrary electromagnetism.
\end{enumerate}
In the language of Newton--Cartan geometry the data $(\mathrm{d}t,\;h,\;A_\mu)$
are the clock form, the spatial metric and the $U(1)$ connection; this is the
statement, due to Son and Wingate (2006) for the unitary Fermi gas and
placed on a general footing by Jensen, that Galilean field theories couple to
Newton--Cartan geometry with a particle-number gauge field, and that
the cleanest way to implement the required boost symmetry is to obtain the
background by null reduction.

%.....................................................................
\subsection{The quantum Hall problem}
\label{sec:nr-qhe}
%.....................................................................

Take $d=2$, $U=0$ and a uniform field $B=\partial_xA_y-\partial_yA_x>0$, for instance
$\mathbf{A}=\tfrac12B(-y,x)$, and set $\hbar=1$. The reduced Hamiltonian is
$H=\boldsymbol{\pi}^2/2m$ with kinetic momentum $\boldsymbol{\pi}=\mathbf{p}-m\mathbf{A}$.

\begin{nrprop}[Landau levels]\label{nr:prop-landau}
The spectrum of $H=\boldsymbol{\pi}^2/2m$ is
\begin{equation}
  E_n=B\bigl(n+\tfrac12\bigr),\qquad n=0,1,2,\dots,
  \label{nr:landau}
\end{equation}
independent of $m$, and each level has $mB/2\pi$ states per unit area.
\end{nrprop}
\begin{proof}
From $[x_i,p_j]=i\delta_{ij}$ one finds $[\pi_x,\pi_y]=imB$. Define
$a=(\pi_x+i\pi_y)/\sqrt{2mB}$. Then $[a,a^\dagger]=1$ and
$a^\dagger a=(\boldsymbol\pi^2-mB)/(2mB)$, so
$H=\boldsymbol\pi^2/2m=B\,(a^\dagger a+\tfrac12)$.
For the degeneracy, define guiding-centre operators
$R_x=x+\pi_y/mB$, $R_y=y-\pi_x/mB$. They commute with $\boldsymbol\pi$
and, hence, with $H$; and $[R_x,R_y]=-i\ell^2$ with $\ell^2=1/mB$. Each state
therefore occupies a phase-space cell $2\pi\ell^2$ in the guiding-centre plane,
which gives $1/2\pi\ell^2=mB/2\pi$ states per unit area.
\end{proof}

The level spacing $B$ does not depend on $m$. This is the equivalence principle
again: the ``charge'' in the Lorentz force is the mass, so the ratio of charge
to mass, which fixes the cyclotron frequency $eB/m$ of a real electron,
equals one for every particle of the geometry. The filling factor of a state
with number density $n$ is $\nu=2\pi n/mB$.

\begin{nrprop}[Kohn's theorem]\label{nr:prop-kohn}
Let $N$ particles with masses $m_a$ and charge equal to mass move in the uniform field $B$ with
$H=\sum_a\boldsymbol\pi_a^2/2m_a+V$, where $V$ depends only on the differences
$\mathbf{x}_a-\mathbf{x}_b$ and is otherwise arbitrary. Then
$H=\boldsymbol\Pi^2/2M+H_{\rm rel}$ with $\boldsymbol\Pi=\sum_a\boldsymbol\pi_a$,
$M=\sum_am_a$ and $[\boldsymbol\Pi,H_{\rm rel}]=0$, so the centre of mass has the
spectrum \eqref{nr:landau} whatever $V$ is. In particular a uniform
oscillating force excites the system only at the frequency $B$.
\end{nrprop}
\begin{proof}
Write $\tilde{\boldsymbol\pi}_a=\boldsymbol\pi_a-(m_a/M)\boldsymbol\Pi$, so that
$\sum_a\tilde{\boldsymbol\pi}_a=0$. Because $[\pi_{b,x},\pi_{a,y}]=im_aB\,\delta_{ab}$,
\[
  [\Pi_x,\tilde\pi_{a,y}]=im_aB-\frac{m_a}{M}\,[\Pi_x,\Pi_y]
  =im_aB-\frac{m_a}{M}\,iMB=0 ,
\]
and likewise for the other components. The cross term in
$\sum_a\boldsymbol\pi_a^2/2m_a$ is proportional to $\boldsymbol\Pi\cdot\sum_a\tilde{\boldsymbol\pi}_a=0$, so
$\sum_a\boldsymbol\pi_a^2/2m_a=\boldsymbol\Pi^2/2M+\sum_a\tilde{\boldsymbol\pi}_a^2/2m_a$.
Also $[\boldsymbol\Pi,\mathbf{x}_a-\mathbf{x}_b]=0$, so $V$ commutes with $\boldsymbol\Pi$. Therefore
$H_{\rm rel}=\sum_a\tilde{\boldsymbol\pi}_a^2/2m_a+V$ commutes with $\boldsymbol\Pi$, and
the centre-of-mass part obeys $[\Pi_x,\Pi_y]=iMB$, which is the Landau problem for mass $M$:
spectrum $B(n+\tfrac12)$. A uniform force couples to
$\sum_am_a\mathbf{x}_a$, which is a linear combination of $\boldsymbol\Pi$ and the
total guiding centre $\sum_am_a\mathbf{R}_a$ (the $\mathbf{R}_a$ defined particle by
particle as in Proposition~\ref{nr:prop-landau}); the latter commutes with $H$,
so only the centre-of-mass oscillator is excited.
\end{proof}

This is the statement that the cyclotron resonance of an interacting electron
gas is not shifted by interactions (Kohn 1961). In the present picture it is
the Galilean boost charge $\mathbf{K}=\sum_am_a\mathbf{x}_a-t\mathbf{P}$ at work:
its central charge is $M$, the total of the charges, which is why the argument
needs the same charge-to-mass ratio for all particles.

\paragraph{Hall conductivity at zero wavenumber.}
Add a uniform acceleration field $\mathbf{a}=-\nabla U$. A classical particle
drifts with the velocity $\mathbf{v}_d$ that cancels the total force, $\mathbf{a}+\mathbf{v}_d\times B\hat{\mathbf z}=0$, i.e.
$\mathbf{v}_d=\mathbf{a}\times\hat{\mathbf z}/B$. The quantum statement is exact
for any Galilean-invariant state. A boost by $\mathbf{v}_d$ is a symmetry of the
ambient space, and it maps the fields $(\mathbf{a},B)$ to $(\mathbf{a}+\mathbf{v}_d\times B\hat{\mathbf z},B)=(0,B)$. In the boosted frame there is
only the magnetic field, and a gapped uniform state carries no current;
back in the original frame the whole fluid moves at $\mathbf{v}_d$. The
number current is therefore
\begin{equation}
  \mathbf{j}=n\,\mathbf{v}_d=\frac{n}{B}\,\mathbf{a}\times\hat{\mathbf z},
  \qquad \sigma_H=\frac{n}{B}=\frac{\nu\,m}{2\pi},
  \label{nr:sigmaH}
\end{equation}
the quantised Hall conductivity. Here $\sigma_H$ relates the number current to the
acceleration field; multiplying by the charge squared (the charge is $e=m$)
gives the electric $\sigma_{xy}=\nu e^2/h$ in the usual units. The value is fixed by the boost
symmetry alone and does not depend on the interactions that produce a gapped state
at a given $\nu$.

\paragraph{Beyond zero wavenumber (results quoted from the literature).}
At finite wavenumber $q$ the boost argument no longer fixes the answer, and
the extra structure is where the Newton--Cartan formalism is used in current
work. Hoyos and Son (2012) and Bradlyn, Goldstein and Read (2012) showed that the
$q^2$ correction to the Hall conductivity has a coefficient related to the
\emph{Hall viscosity} $\eta_H$, the non-dissipative part of the stress response
to shear. For quantum Hall states $\eta_H=\tfrac12\hbar\,n\,\bar s$, with $\bar s$ the
mean orbital spin per particle (Read 2009), and the combination $2\bar s$ is
the topological ``shift'' of the state on a sphere. Son (2013) proposed that the
lowest Landau level, obtained as the limit $m\to0$ at fixed charge with
$g$-factor $2$ (a Zeeman-like term then puts the lowest Landau level at exactly zero energy), is described by an effective field theory
built on Newton--Cartan geometry. The $m\to0$ limit is smooth in
this formulation. Geracie, Son, Wu and Wu (2015) then used the resulting
diffeomorphism Ward identities to relate the leading wavenumber dependence of
the Hall conductivity to the shift, to derive sum rules for the stress-tensor
spectral densities, and to obtain an inequality between the $k^4$ coefficient
of the projected static structure factor and the shift. We do not reproduce
these derivations here; the point for this chapter is that the symmetry
structure they exploit is the one that the null reduction makes manifest.

%.....................................................................
\subsection{Schr\"odinger spacetime and non-relativistic holography}
\label{sec:nr-holography}
%.....................................................................

In the holographic correspondence a quantum field theory is described by gravity in
a space with one more dimension. For theories with Schr\"odinger symmetry, the
symmetry of the free Schr\"odinger equation and of many strongly interacting
systems such as cold atoms at unitarity (Son and Wingate 2006; Nishida and Son
2007), a gravity dual is obtained by a null construction. With
$\mathbf{x}\in\mathbb{R}^d$, a radial coordinate $z>0$, and our $(t,s)$,
\begin{equation}
  g_{\rm Sch}=\frac{L^2}{z^2}\Bigl(\mathrm{d}\mathbf{x}^2+\mathrm{d}z^2
  +2\,\mathrm{d}t\,\mathrm{d}s-\frac{\mathrm{d}t^2}{z^2}\Bigr).
  \label{nr:sch}
\end{equation}
(In the literature $s=-\xi$ and the metric is written
$L^2[-\mathrm{d}t^2/z^4+(-2\,\mathrm{d}t\,\mathrm{d}\xi+\mathrm{d}\mathbf{x}^2+\mathrm{d}z^2)/z^2]$; a
coefficient $\beta^2$ in front of $\mathrm{d}t^2/z^4$ can be set to $1$ by
rescaling $(t,s)\mapsto(at,s/a)$.) The metric has $d+3$ dimensions.

\begin{nrprop}[Structure and symmetries of Schr\"odinger spacetime]\label{nr:prop-sch}
Let $g_{\rm Sch}$ be the metric \eqref{nr:sch} on $\tilde M_{\rm Sch}=\{(\mathbf{x},t,s,z):z>0\}$.
\begin{enumerate}
\item The bracket in \eqref{nr:sch} is the Eisenhart--Duval metric
\eqref{nr:gPhi}, with $d+1$ space dimensions (the coordinates $\mathbf{x}$ and $z$)
and the potential $\Phi=1/(2z^2)$; the full metric \eqref{nr:sch} is
$L^2/z^2$ times it. In particular $\xi=\partial_s$ is a null Killing vector
and the quotient by $\xi$ carries a Galilean structure.
\item The following maps are isometries (for every $d$; checked by direct
substitution):
\begin{align*}
  \text{boosts:}\quad & \mathbf{x}\mapsto\mathbf{x}+\mathbf{v}t,\quad
     s\mapsto s-\mathbf{v}\cdot\mathbf{x}-\tfrac12v^2t,\\
  \text{dilatations:}\quad & (t,\mathbf{x},z,s)\mapsto(\lambda^2t,\lambda\mathbf{x},\lambda z,s),\\
  \text{special conformal:}\quad & (t,\mathbf{x},z)\mapsto\frac{(t,\mathbf{x},z)}{1+ct},\\
     & s\mapsto s+\frac{c}{2}\,\frac{\mathbf{x}^2+z^2}{1+ct}.
\end{align*}
Together with translations of $t,\mathbf{x},s$ and rotations of $\mathbf{x}$ they generate
the Schr\"odinger group, of dimension $\tfrac12d(d-1)+2d+4$ ($13$ for $d=3$).
\end{enumerate}
\end{nrprop}
\begin{proof}
(1) The factor in brackets is $\mathrm{d}\mathbf{x}^2+\mathrm{d}z^2+2\,\mathrm{d}t\,\mathrm{d}s-2\Phi\,\mathrm{d}t^2$
with $-2\Phi=-1/z^2$. (2) Boosts do not change $z$ or $t$ and preserve
$\mathrm{d}\mathbf{x}^2+2\,\mathrm{d}t\,\mathrm{d}s$ by the computation after
\eqref{nr:finite}. Under dilatation, $\mathrm{d}\mathbf{x}^2/z^2$, $\mathrm{d}z^2/z^2$, $\mathrm{d}t\,\mathrm{d}s/z^2$
and $\mathrm{d}t^2/z^4$ each scale by $\lambda^0$. The special conformal case is a
direct, if longer, substitution.
\end{proof}

The dynamical exponent $2$ (time scales as the square of length) is built
into the dilatation, as in a free Schr\"odinger particle. The Schr\"odinger group is
the centraliser of $\partial_s$ in the conformal group of the flat ambient
space $\mathbb{R}^{d+1,1}$ of this chapter (Burdet, Perrin and Sorba 1973; Duval and
Horv\'athy 2009), which is why Schr\"odinger symmetry and the null reduction
go together. The boundary $z\to0$ carries the coordinates $(t,\mathbf{x},s)$ and the
boundary theory is a Galilean field theory in $d$ space dimensions whose particle
number is the momentum along $s$.

\begin{nrprop}[A scalar field in Schr\"odinger spacetime]\label{nr:prop-schscalar}
A scalar of bulk mass $\mu$, $(\Box-\mu^2)\Phi=0$, with dependence
$\Phi=e^{-i\omega t+ims+i\mathbf{k}\cdot\mathbf{x}}\phi(z)$ obeys
\begin{equation}
  z^2\phi''-(d+1)\,z\,\phi'-\bigl(q^2z^2+m^2+\mu^2L^2\bigr)\phi=0,
  \qquad q^2:=\mathbf{k}^2-2m\omega .
  \label{nr:radial}
\end{equation}
Its solution that decays into the bulk is $\phi=z^{(d+2)/2}K_\nu(qz)$ with
\begin{equation}
  \nu=\sqrt{\tfrac14(d+2)^2+m^2+\mu^2L^2},
  \label{nr:nu}
\end{equation}
which is oscillatory for $q^2<0$.
\end{nrprop}
\begin{proof}
The inverse metric is $g^{ij}=g^{zz}=(z^2/L^2)\delta^{ij}$, $g^{ts}=z^2/L^2$,
$g^{ss}=1/L^2$, and $\sqrt{-g}=(L/z)^{d+3}$. In
$\Box\Phi=(\sqrt{-g})^{-1}\partial_\mu(\sqrt{-g}\,g^{\mu\nu}\partial_\nu\Phi)$
the $z$-derivatives give $(z^2/L^2)\bigl[\phi''-(d+1)\phi'/z\bigr]$, the $\mathbf{x}$-derivatives
give $-(z^2/L^2)\mathbf{k}^2\phi$, the mixed term $2g^{ts}\partial_t\partial_s$ gives
$+2(z^2/L^2)\,m\omega\,\phi$, and $g^{ss}\partial_s^2$ gives $-(m^2/L^2)\phi$.
Multiplying $(\Box-\mu^2)\Phi=0$ by $L^2$ gives \eqref{nr:radial}. With
$\phi=z^{a}f$, $a=\tfrac12(d+2)$, it becomes the modified Bessel equation of order
$\nu$ with $\nu^2=a^2+m^2+\mu^2L^2$.
\end{proof}

Three things can be read from \eqref{nr:radial}--\eqref{nr:nu}.
\begin{enumerate}
\item \emph{The Galilean dispersion is the bulk light cone.} The solution changes
character at $q^2=0$, i.e.\ at $\omega=\mathbf{k}^2/2m$. This is the null
condition $\mathbf{p}^2+2p_tp_s=0$ of the ambient space, with $p_t=-\omega$ and
$p_s=m$ (Section~\ref{sec:nr-dynamics}): a free Galilean particle of mass $m$ is a null
ray of the ambient geometry, and the continuum of the boundary two-point
function begins on that dispersion curve.
\item \emph{Correlators and dimensions.} The standard holographic prescription
(compare the coefficients of $z^{\frac12(d+2)\mp\nu}$ in $K_\nu$) gives the
retarded two-point function of the dual operator, for non-integer $\nu$, as
$G(\omega,\mathbf{k})\propto(\mathbf{k}^2-2m\omega)^{\nu}$ up to contact terms,
with scaling dimension
$\Delta=\tfrac12(d+2)+\nu$, in the sense that
$G$ scales with weight $2\Delta-(d+2)=2\nu$ under the dilatation with $[t]=-2$,
$[\mathbf{x}]=-1$. The dimension depends on $m$ through the $\mathrm{d}t^2/z^4$ term:
operators carrying different particle number have different dimensions. In
the boundary theory the dimension of an operator equals the energy of the
corresponding state in a harmonic trap, in units of the trap frequency (the
operator--state correspondence of Nishida and Son 2007).
\item \emph{Superselection and particle number.} Each value of the momentum
$m$ along $s$ gives an independent sector, the holographic form of Bargmann's mass
superselection rule. If $s$ is made periodic, $m$ is quantised, as in
Section~\ref{sec:nr-quantum}, and $m$ counts particles.
\end{enumerate}

\begin{nrrem}[Status of the construction]
The metric \eqref{nr:sch} is not a vacuum solution of Einstein's equations in
$d+3$ dimensions: it requires matter, for example a massive vector field
along $\mathrm{d}t$ (Son 2008; Balasubramanian and McGreevy 2008). It should be
regarded as a solvable model with the right symmetries, not as the known dual of
a particular atomic system. Its use has been to compute
transport and thermodynamic properties of non-relativistic conformal
systems at strong coupling and to explore the geometry (Newton--Cartan
boundary, torsion, related Lifshitz duals obtained by further reduction; see
the review by Taylor 2016). In this application the null direction $s$ is still the
particle-number direction of the chapter, and the new, emergent
direction is the radial coordinate $z$. The distinction between the two is
what the observability discussion of the preceding section would
lead one to expect.
\end{nrrem}

%---------------------------------------------------------------------
\section{Summary}
\label{sec:nr-summary}
%---------------------------------------------------------------------

\begin{enumerate}
\item The free Galilean particle has a Lagrangian that is quasi-invariant under
boosts, $\delta L=\mathrm{d}F/\mathrm{d}t$ with $F\propto m$; the Noether
charges therefore close on a central extension, $\{K_i,P_j\}=m\delta_{ij}$
\eqref{nr:KP}. The extension is cohomologically non-trivial
(Proposition~\ref{nr:prop-nontrivial}).
\item Galilean spacetime is the orbit space of the null parallel vector field
$\xi=\partial_s$ on $\mathbb{R}^{d+1,1}$ with $g=\mathrm{d}\mathbf{x}^2+2\,\mathrm{d}t\,\mathrm{d}s$. The
quotient inherits a clock one-form and a degenerate spatial metric
(Proposition~\ref{nr:prop-induced}) and nothing else.
\item Mass is $p_s$, conserved because $\xi$ is Killing. Free null geodesics
project to uniform motion with $E=\mathbf{p}^2/2m$ \eqref{nr:dispersion}, and
$S=-m\,\Delta s$ \eqref{nr:action}.
\item The isometries of $\tilde M$ commuting with $\xi$ form the Bargmann group,
of dimension $\tfrac12d(d-1)+2d+2$ \eqref{nr:dim}. The cocycle is realised as a shift
of $s$ in the group commutator of a boost and a translation \eqref{nr:groupcomm}.
\item The deformation $-2\Phi\,\mathrm{d}t^2$ of the metric produces Newtonian
gravity with universal free fall \eqref{nr:newton}, and $R_{\mu\nu}=\nabla^2\Phi\,\tau_\mu\tau_\nu$
\eqref{nr:Ricci}. Quantum mechanically the same structure gives the
Schr\"odinger equation and the boost phase \eqref{nr:boostphase}.
\item The Bargmann Casimir $2MH-\mathbf{P}^2$ is the squared mass of the
ambient Poincar\'e group \eqref{nr:twoCasimirs}; the central charge of the
Galilean algebra is the light-front momentum $P^+$.
\item The reduction is an exact projection with no limit involved, and differs
from the $c\to\infty$ contraction (Table~\ref{tab:nr-routes}).
\item The null reduction has a spacelike counterpart in Klein's and Rumer's
five-dimensional optics: the same eikonal construction in the same flat space with a
spacelike extra direction gives relativistic mechanics, and in both cases the action is
$-m$ times the displacement along the extra coordinate
(Table~\ref{tab:nr-rumer}, Section~\ref{sec:nr-rumer}).
\item In applications, a Galilean field theory on a background is a massless field
on a Bargmann space: mass becomes a $U(1)$ charge and boosts come for free
(Proposition~\ref{nr:prop-reducedaction}); the quantum Hall problem has $m$-independent
Landau levels and an interaction-independent cyclotron resonance
(Propositions~\ref{nr:prop-landau}, \ref{nr:prop-kohn}); and non-relativistic holography
uses Schr\"odinger spacetime \eqref{nr:sch}, whose boundary dispersion is the ambient null
cone (Proposition~\ref{nr:prop-schscalar}).
\end{enumerate}

%---------------------------------------------------------------------
\section{Problems}
%---------------------------------------------------------------------

\begin{enumerate}
\item \textbf{Oscillator.} Take $d=1$ and $\Phi=\tfrac12\omega^2x^2$ in
\eqref{nr:gPhi}. Show that the projected null geodesics satisfy
$\ddot x=-\omega^2x$ with no dependence on $m$. For $x(t)=A\cos\omega t$ show
that
\[
  s(t)=s_0+\tfrac14 A^2\omega\,\sin 2\omega t .
\]
Deduce that $s$ returns to its initial value after one period, and that the
time average of the Lagrangian over a period vanishes. (This is the virial
theorem for the harmonic oscillator, $\langle T\rangle=\langle V\rangle$,
read as periodicity of $s$.)

\item \textbf{Many particles.} Let $N$ particles be described by coordinates
$(\mathbf{x}_a,s_a)$, $a=1,\dots,N$, with a common time $t$, and let the
Bargmann group act diagonally, $\mathbf{x}_a\mapsto\mathbf{x}_a+\mathbf{v}t$,
$s_a\mapsto s_a-\mathbf{v}\cdot\mathbf{x}_a-\tfrac12v^2t$. Show that the total
charges are $K_i=\sum_a(x_{a,i}\,p_{s_a}-t\,p_{a,i})$, $P_i=\sum_a p_{a,i}$,
$M=\sum_a p_{s_a}$, and that $\{K_i,P_j\}=M\delta_{ij}$. Conclude that the mass
of the composite is the sum of the masses, independently of interactions
depending on relative positions.

\item \textbf{Boosts and time translations.} Let $T_b$ be the translation
$t\mapsto t+b$ and $B_{\mathbf v}$ the boost with $R=1$ in \eqref{nr:finite}.
Show that
\[
  B_{\mathbf v}\,T_b=T_{\mathbf{v}b}\;Z\bigl(-\tfrac12v^2b\bigr)\;T_b\,B_{\mathbf v},
\]
where $T_{\mathbf a}$ is a spatial translation and $Z(c)$ the shift of $s$ by $c$.
On the base $M$ the commutator of a boost and a time translation is a spatial
translation by $\mathbf{v}b$; on $\tilde M$ there is, in addition, a shift of
$s$ by $-\tfrac12 v^2b$. Compare the latter with the term $\tfrac12mu^2t$ in
$F_{\mathbf u}$ of \eqref{nr:quasi}.

\item \textbf{Spin Casimir.} In $d=3$, with the Poisson algebra
\eqref{nr:KP}--\eqref{nr:KK} and the rotation relations
$\{J_i,K_j\}=\epsilon_{ijk}K_k$, $\{J_i,P_j\}=\epsilon_{ijk}P_k$ (where
$J_k=\tfrac12\epsilon_{kij}J_{ij}$), verify that
$\mathbf{S}=\mathbf{J}-M^{-1}\,\mathbf{K}\times\mathbf{P}$ satisfies
$\{K_i,S_j\}=0$ and $\{P_i,S_j\}=0$, hence that $\mathbf{S}^2$ is a Casimir
function of the Bargmann algebra.

\item \textbf{Spacelike reduction.} In $\mathbb{R}^{d+1,1}$ with
$g=-\mathrm{d}t^2+\mathrm{d}\mathbf{x}^2+\mathrm{d}y^2$, show that null geodesics with
$p_y=m$ have $\mathrm{d}y=\mathrm{d}\tau$ and $E=\sqrt{\mathbf{p}^2+m^2}$, and verify
\eqref{nr:restsub}. Show that the null coordinate $s=(y-t)/\sqrt2$ together with
$t'=(y+t)/\sqrt2$ turns $g$ into \eqref{nr:metric}, and find the relation between the
spacelike and null momenta, $p_s=(p_y+E)/\sqrt2$ and $E'=-p_{t'}=(E-p_y)/\sqrt2$.
For slow motion with $p_y=m$, show that $E'\approx\mathbf{p}^2/(2\sqrt2\,m)=\mathbf{p}^2/2p_s$ with
$p_s\approx\sqrt2\,m$.

\item \textbf{Schr\"odinger isometries.} For $d=1$, verify that the special conformal map
of Proposition~\ref{nr:prop-sch} preserves \eqref{nr:sch}, and that the opposite sign
of the shift in $s$ does not. Compute the Killing vectors that generate the
three families, and check the commutator of the dilatation generator with the
generator of time translations.

\item \textbf{Radial equation.} Derive \eqref{nr:radial} for $d=2$. Show that for
$q^2>0$ the decaying solution behaves as $z^{\frac12(d+2)-\nu}$ near $z=0$ with
a coefficient of $z^{\frac12(d+2)+\nu}$ proportional to $q^{2\nu}$, and explain why this
gives $G\propto(\mathbf{k}^2-2m\omega)^\nu$.

\item \textbf{Landau problem with an acceleration field.} In the setting of
Proposition~\ref{nr:prop-landau}, add $U=-\mathbf{a}\cdot\mathbf{x}$ and show that, after
shifting $\boldsymbol\pi$ by a constant,
$H=B(a^\dagger a+\tfrac12)-\tfrac12m\,v_d^2-m\,\mathbf{a}\cdot\mathbf{R}$ with $v_d=a/B$, and
that the guiding centre drifts at $\mathbf{v}_d=\mathbf{a}\times\hat{\mathbf z}/B$.
\end{enumerate}

%---------------------------------------------------------------------
\section*{Notes and references}
%---------------------------------------------------------------------

The classification of ray representations of the Galilei group, the
central extension and the mass superselection rule are due to
Bargmann~(1954) and L\'evy-Leblond~(1963); the geometric formulation in terms of
moment maps and cocycles is Souriau's~(1970). The use of an extra dimension to
turn Newtonian trajectories into geodesics goes back to Eisenhart~(1929). Its
relation to Newton--Cartan theory (Cartan 1923/24; K\"unzle 1972) via
\emph{Bargmann structures} was established by Duval, Burdet, K\"unzle and
Perrin~(1985), and developed by Duval, Gibbons and Horv\'athy~(1991). The
light-front viewpoint is Dirac's~(1949), and null identifications are the basis
of discrete light-cone quantisation (see Seiberg 1997). Klein's five-dimensional wave equation and Rumer's
5-optics, which treats the action as a coordinate, are the spacelike precursor of
this construction (Section~\ref{sec:nr-rumer}). The applications of
Section~\ref{sec:nr-applications} follow Son and Wingate (2006), Jensen (arXiv:1408.6855),
Son (2008, 2013), Balasubramanian and McGreevy (2008) and Geracie, Son, Wu and Wu
(2015); the derivations given there of Propositions~\ref{nr:prop-reducedaction}--\ref{nr:prop-schscalar}
are elementary and are written out for completeness.

\begin{itemize}
\item V.~Bargmann, On unitary ray representations of continuous groups,
      \emph{Ann.\ Math.} \textbf{59} (1954) 1.
\item J.-M.~L\'evy-Leblond, Galilei group and nonrelativistic quantum mechanics,
      \emph{J.\ Math.\ Phys.} \textbf{4} (1963) 776.
\item J.-M.~L\'evy-Leblond, Galilei group and Galilean invariance, in
      \emph{Group Theory and Its Applications}, Vol.~II, ed.\ E.~M.~Loebl
      (Academic Press, 1971).
\item J.-M.~Souriau, \emph{Structure des syst\`emes dynamiques} (Dunod, 1970).
\item L.~P.~Eisenhart, Dynamical trajectories and geodesics,
      \emph{Ann.\ Math.} \textbf{30} (1929) 591.
\item \'E.~Cartan, Sur les vari\'et\'es \`a connexion affine et la th\'eorie de la
      relativit\'e g\'en\'eralis\'ee, \emph{Ann.\ \'Ec.\ Norm.\ Sup.}
      \textbf{40} (1923) 325 and \textbf{41} (1924) 1.
\item H.~P.~K\"unzle, Galilei and Lorentz structures on space-time:
      comparison of the corresponding geometry and physics,
      \emph{Ann.\ Inst.\ H.\ Poincar\'e A} \textbf{17} (1972) 337.
\item C.~Duval, G.~Burdet, H.~P.~K\"unzle and M.~Perrin, Bargmann structures and
      Newton--Cartan theory, \emph{Phys.\ Rev.\ D} \textbf{31} (1985) 1841.
\item C.~Duval, G.~W.~Gibbons and P.~Horv\'athy, Celestial mechanics, conformal
      structures, and gravitational waves, \emph{Phys.\ Rev.\ D} \textbf{43}
      (1991) 3907.
\item E.~In\"on\"u and E.~P.~Wigner, On the contraction of groups and their
      representations, \emph{Proc.\ Natl.\ Acad.\ Sci.\ USA} \textbf{39} (1953) 510.
\item H.~Bacry and J.-M.~L\'evy-Leblond, Possible kinematics,
      \emph{J.\ Math.\ Phys.} \textbf{9} (1968) 1605.
\item P.~A.~M.~Dirac, Forms of relativistic dynamics,
      \emph{Rev.\ Mod.\ Phys.} \textbf{21} (1949) 392.
\item T.~Kaluza, Zum Unit\"atsproblem der Physik, \emph{Sitzungsber.\ Preuss.\ Akad.\ Wiss.}
      (1921) 966.
\item O.~Klein, Quantentheorie und f\"unfdimensionale Relativit\"atstheorie,
      \emph{Z.\ Phys.} \textbf{37} (1926) 895.
\item Yu.~B.~Rumer, Action as a space coordinate (in Russian),
      \emph{Zh.\ Eksp.\ Teor.\ Fiz.} \textbf{19} (1949); and
      \emph{Issledovaniya po 5-optike} (Studies on 5-optics) (Gostekhizdat,
      Moscow, 1956).
\item N.~Redington, On the significance of the fifth coordinate in Wesson's
      version of Kaluza--Klein theory, arXiv:gr-qc/9701062 (1997).
\item START in a five-dimensional conformal domain (on Keller's START theory and
      Rumer's 5-optics), arXiv:1111.5540.
\item N.~Seiberg, Why is the Matrix model correct?,
      \emph{Phys.\ Rev.\ Lett.} \textbf{79} (1997) 3577.
\item D.~T.~Son and M.~Wingate, General coordinate invariance and conformal
      invariance in nonrelativistic physics: unitary Fermi gas,
      \emph{Ann.\ Phys.} \textbf{321} (2006) 197.
\item Y.~Nishida and D.~T.~Son, Nonrelativistic conformal field theories,
      \emph{Phys.\ Rev.\ D} \textbf{76} (2007) 086004.
\item D.~T.~Son, Toward an AdS/cold atoms correspondence: a geometric
      realization of the Schr\"odinger symmetry, \emph{Phys.\ Rev.\ D}
      \textbf{78} (2008) 046003.
\item K.~Balasubramanian and J.~McGreevy, Gravity duals for non-relativistic CFTs,
      \emph{Phys.\ Rev.\ Lett.} \textbf{101} (2008) 061601.
\item G.~Burdet, M.~Perrin and P.~Sorba, About the non-relativistic structure of
      the conformal algebra, \emph{Commun.\ Math.\ Phys.} \textbf{34} (1973) 85.
\item C.~Duval and P.~A.~Horv\'athy, Non-relativistic conformal symmetries and
      Newton--Cartan structures, \emph{J.\ Phys.\ A} \textbf{42} (2009) 465206.
\item K.~Jensen, On the coupling of Galilean-invariant field theories to curved
      spacetime, arXiv:1408.6855 (published in \emph{SciPost Phys.}, 2018).
\item W.~Kohn, Cyclotron resonance and de Haas--van Alphen oscillations of an
      interacting electron gas, \emph{Phys.\ Rev.} \textbf{123} (1961) 1242.
\item C.~Hoyos and D.~T.~Son, Hall viscosity and electromagnetic response,
      \emph{Phys.\ Rev.\ Lett.} \textbf{108} (2012) 066805.
\item B.~Bradlyn, M.~Goldstein and N.~Read, Kubo formulas for viscosity: Hall
      viscosity, Ward identities, and the relation with conductivity,
      \emph{Phys.\ Rev.\ B} \textbf{86} (2012) 245309.
\item N.~Read, Non-Abelian adiabatic statistics and Hall viscosity in quantum Hall
      states and $p_x+ip_y$ paired superfluids, \emph{Phys.\ Rev.\ B}
      \textbf{79} (2009) 045308.
\item D.~T.~Son, Newton--Cartan geometry and the quantum Hall effect,
      arXiv:1306.0638.
\item M.~Geracie, D.~T.~Son, C.~Wu and S.-F.~Wu, Spacetime symmetries of the
      quantum Hall effect, \emph{Phys.\ Rev.\ D} \textbf{91} (2015) 045030.
\item M.~Taylor, Lifshitz holography, arXiv:1512.03554.
\end{itemize}
